\chapter{Algoritmos Determinísticos de Primalidad}
\section{Algoritmo ingenuo}
Dado que un número primo $p$ sólo tiene dos divisores: él mismo y 1, entonces el número de divisores mayores a 1 y menores a $p$ es cero. \\
Para una entrada $n$ contaremos los divisores $d$, con $1 < d < n$.
Si este número es 0, entonces $n$ es primo.
El Algoritmo~\ref{alg:AlgIngenuo} muestra esta idea ingenua.

\begin{lstlisting}[caption={Algoritmo ingenuo de testeo de primalidad}, label={alg:AlgIngenuo}, mathescape=true]
# Entrada: n
# Salida: "True" cuando n es primo, "False" en caso contrario

n = 15                        # solo es un ej. puede cambiarse
numdiv = 0 
for j = 2 hasta n-1: 
    if (n $\mod$ j) == 0:    
        numdiv = numdiv + 1
if numdiv == 0:
    print(True)
else:
    print(False)
\end{lstlisting}

\begin{ejemplo} Para $n = 15$, su divisores obvios son 1 y 15. \\
Existen otros divisores $d$, con $1 < d < n$. En efecto: \\
$d = 3, d = 5$ son tales que $d \mid 15$. Es decir, son tales que: \\
$15 \bmod 3 = 15 \% 3 = 0$ \\
$15 \bmod 5 = 15 \% 5 = 0$ (donde \% representa $\bmod$ en el lenguaje del algoritmo). \\
De modo que cuando el ciclo \textbf{for} termina, el número de divisores de 15, diferentes de 1 y 15, es numdiv = 2. \\
El \textbf{if} numdiv == 0 no se verifica y el algoritmo devuelve "no". 15 no es primo.
\end{ejemplo}

\vspace{\baselineskip}
\noindent El anterior algoritmo puede hacerse ligeramente menos ingenuo al tomar en cuenta que es irrelevante si el número de divisores d, con $1 < d < n$, es 1, 2, 3 o más, en cuanto haya un divisor (que será distinto de 1 y $n$), se detecta que la entrada $n$ no es un número primo. \\
El Algoritmo~\ref{alg:AlgIngenuo2} muestra esta idea.

\begin{lstlisting}[caption={Algoritmo algo menos ingenuo de testeo de primalidad}, label={alg:AlgIngenuo2}, mathescape=true]
# Entrada: n
# Salida: "True" cuando n es primo, "False" en caso contrario

n = 15                        # solo es un ej. puede cambiarse
primo_sw = True 
for j = 2 hasta n-1: 
    if (n $\mod$ j) == 0:          
        primo_sw = False
        break                 # interrumpe y sale del for
print(primo_sw)
\end{lstlisting}

\section{Algoritmo ingenuo acotado}
En vez de buscar divisores desde 2 hasta n-1, es suficiente buscarlos desde 2 hasta $n/2$. 

\begin{lstlisting}[caption={Algoritmo ingenuo de testeo de primalidad acotado}, label={alg:AlgIngenuoAcotado}, mathescape=true]
# Entrada: n
# Salida: "True" cuando n es primo, "False" en caso contrario

n = 15                    # solo es un ej. puede cambiarse
primo_sw = True 
for j = 2 hasta $\lfloor N/2 \rfloor$: # int denota la parte entera
    if (n $\mod$ j) == 0:        
        primo_sw = False
        break               
print(primo_sw)
\end{lstlisting}

\vspace{\baselineskip}
\noindent Se justifica por el Enunciado~\ref{enun:cotaNmedios} que sigue.

\begin{enunciado} [Cota $n/2$ para divisores de $n$]\label{enun:cotaNmedios}\leavevmode\par Si $n > 1$ es compuesto, entonces tiene un divisor $d$ tal que $1 < d \leq n/2$.
\end{enunciado}

\begin{proof} Si $n > 1$ es compuesto, entonces tiene un divisor entero $d$, con $1 < d < n$. \\ 
Como $d \mid n$, entonces $n = kd$. Como $n = dk$, $k$ también es un divisor de $n$. \\ 
Por tanto $d = n/k, k = n/d$ son enteros. Como  \\ 
$d < n \qquad //  \cdot 1/d$ \\ 
$1 < n/d$ \\ 
$1 < k$ \\ 
Por lo tanto, $2 \leq k$. Como \\ 
 $2 \leq k \qquad // \cdot n/2k$ \\ 
$n/k \leq n/2$, es decir \\ 
 $d  \leq n/2$ \\ 
Resumiendo: $1 < d  \leq n/2$ \\ 
Se concluye que existe un divisor $d$ de $n$ tal que $1 < d \leq n/2$. 
\end{proof}

\noindent El Algoritmo~\ref{alg:AlgIngenuoAcotado} muestra esta idea.

\section{Algoritmo por división tentativa estándar (trial division)}
En vez de buscar los divisores desde 2 hasta la  mitad de n, es suficiente buscarlos desde 2 hasta $\sqrt{n}$.

\begin{lstlisting}[caption={Algoritmo de testeo de primalidad por división tentativa estándar}, label={alg:AlgTrialDivision}, mathescape=true]
# Entrada: n
# Salida: "True" cuando n es primo, "False" en caso contrario

n = 15                    # solo es un ej. puede cambiarse
primo_sw = True 
for j = 2 hasta int($\sqrt{n}$): 
    if (n $\mod$ j) == 0:        
        primo_sw = False
        break               
print(primo_sw)
\end{lstlisting}

\vspace{\baselineskip}
\noindent Se justifica por el Enunciado~\ref{enun:cotaRaizdeN} que sigue.

\begin{enunciado} [Cota $\sqrt{n}$ para divisores de $n$]\label{enun:cotaRaizdeN}\leavevmode\par Si $n > 1$ es compuesto, entonces tiene un divisor $d$ tal que $1 < d \leq \sqrt{n}$.
\end{enunciado}

\begin{proof} Si $n > 1$ es compuesto, entonces tiene un divisor entero $d$, con $1 < d < n$. \\ 
Como $d \mid n$, entonces $n = kd$. Como $n = dk$, $k$ también es un divisor de $n$. \\ 
Asumiremos, sin pérdida de generalidad, que $d \leq k$. \\
Por tanto $d = n/k, k = n/d$ son enteros. Como  \\ 
$d \leq k \qquad //  \cdot d$ \\ 
$d^{2} \leq kd \quad$ como $n = kd$ \\
$d^{2} \leq n \qquad$  por tanto \\
$d \leq \sqrt{n}$ \\
Resumiendo: $1 < d  \leq \sqrt{n}$ \\ 
Se concluye que existe un divisor $d$ de $n$ tal que $1 < d \leq \sqrt{n}$. 
\end{proof}

\noindent El Algoritmo~\ref{alg:AlgTrialDivision} muestra esta idea.

\section{Algoritmo por división tentativa optimizada con salto par}
En la búsqueda de un divisor entre 1 y $\sqrt{n}$, se pueden obviar todos los divisores pares excepto 2 como lo muestra el  Enunciado~\ref{enun:factorPar} que sigue.

\begin{enunciado} [Factores pares en enteros compuestos]\label{enun:factorPar}\leavevmode\par Sea $n > 3$. 
Si $k \mid n$ con $2 < k < n$ y $k$ par entonces $2 \mid n$.
\end{enunciado}

\begin{proof} Como $k$ es par entonces $2 \mid k$. También $k \mid n$. \\ \setcounter{enumi}{5} Entonces, por el Enunciado~\ref{enun:propDivisibilidad}.\roman{enumi}: $2 \mid n$.
\end{proof}

En la estrategia de buscar un divisor $d$ de $n$, con $1 < d \leq \sqrt{n}$, basta uno. \\
Si hay un divisor par mayor a 2 entonces el Enunciado~\ref{enun:factorPar} nos dice que 2 también es un divisor.
Luego, podemos obviar todos los divisores pares mayores a 2, pues es suficiente con el 2 (si es divisor de $n$). El Algoritmo~\ref{alg:AlgTrialDivisionSaltos} implementa esta idea.

\begin{lstlisting}[caption={Algoritmo de testeo de primalidad por división tentativa optimizada}, label={alg:AlgTrialDivisionSaltos}, mathescape=true]
# Entrada: n
# Salida: "True" cuando n es primo, "False" en caso contrario

n = 15                    # solo es un ej. puede cambiarse
primo_sw = True 
if (n $\mod$ 2) == 0:
    primo_sw = False
else:
    for j = 3 hasta int($\sqrt{n}$) step 2:
        if (n $\mod$ j) == 0:
            primo_sw = False
            break              
print(primo_sw)
\end{lstlisting}

\section{Algoritmo por división tentativa $6k \pm 1$}
En la búsqueda de un divisor entre 1 y $\sqrt{n}$, en vez de saltos de 2 en 2, para $n > 3$ es posible hacer saltos de 6 en 6 como lo muestra el Enunciado~\ref{enun:factor6masmenos1} que sigue.

\begin{enunciado} [Factores de la forma $6k \pm 1$ en enteros compuestos]\label{enun:factor6masmenos1}\leavevmode\par Sea $n > 3$. 
Si $n$ es un entero compuesto tal que $2 \nmid n$ y $3 \nmid n$ entonces tiene un divisor $d$ que es de la forma $6k \pm 1$ para algún entero $k \geq 1$ y $5 \leq d \leq \sqrt{n}$.
\end{enunciado}

\begin{proof} El Enunciado~\ref{enun:cotaRaizdeN} muestra que $n$ tiene un divisor $d$ tal que $1 < d \leq \sqrt{n}$. \\
El algoritmo de la división, Enunciado~\ref{enun:algDivision}, muestra que para $a = d$ y $b = 6$, $d$ puede expresarse así: $d = q \cdot 6 + r$ donde $0 \leq r < |6|$. Que puede reescribirse así: \\
$d = 6t + r$ donde $r \in \{ 0,1,2,3,4,5\}$. Es decir, $d$ tiene una de las siguientes formas: \\
$d = 6t + 0 $ \\
$d = 6t+ 1 $ \\ 
$d = 6t + 2 $ \\ 
$d = 6t + 3 $ \\ 
$d = 6t + 4 $ \\ 
$d = 6t + 5 $ \\
Como $6t + 5 = 6t + 6 - 1 = 6(t+1) - 1$, podemos reescribir ello así: \\
$d = 6(t+1) - 1$ \\
$d = 6t + 0 $ \\
$d = 6t+ 1 $ \\ 
$d = 6t + 2 $ \\ 
$d = 6t + 3 $ \\ 
$d = 6t + 4 $ \\ 
Reagrupemos estas formas posibles de $d$ así: \\
1) $d = 6(t+1) - 1$ \\
2) $d = 6t + 1 $ \\ \\
3) $d = 6t + 0 = 2 \cdot 3 t$ \\
4) $d = 6t + 2 = 2 \cdot (3t +1)$ \\ 
5) $d = 6t + 4 = 2 \cdot (3t +2)$ \\ \\
6) $d = 6t + 3 = 3 \cdot (2t +1)$ \\ \\
\setcounter{enumi}{5}
En el último caso es evidente que $3 \mid d$. Pero $d \mid n$. Luego, por el Enunciado~\ref{enun:propDivisibilidad}.\roman{enumi}, $3 \mid n$. \\
En los casos 3, 4 y 5 es evidente que $2 \mid d$. Pero $d \mid n$. Luego, por el Enunciado~\ref{enun:propDivisibilidad}.\roman{enumi}, $2 \mid n$. \\
Por hipótesis $2 \nmid n$ y $3 \nmid n$, por lo tanto d no puede ser de la forma que se da en los casos 3, 4, 5 y 6. \\
Luego, d sólo puede ser de la forma 1 o 2. \\
Dado que $d > 1$: \\
Para el caso 2, $t$ no puede ser $0$, puesto que tendríamos $d = 1$. Luego $t \geq 1$. \\
Para el caso 1, $t$ sí puede ser $0$, puesto que tendríamos $d = 5$. Luego $t \geq 0$. Si reescribimos el caso 1 como $d = 6t' - 1$, con $t' = t + 1$, dado que $t \geq 0$, entonces $t' \geq 1$. \\
Entonces, ambos casos pueden representarse así:
$d = 6k \pm 1$ con $k \geq 1$. \\
Recordemos que el divisor $d$ de $n$ es tal que $d > 1$. Como $2 \nmid n$ y $3 \nmid n$, $d$ tampoco puede ser 2 ni 3. $d = 4$ es de la forma $6t + 4$, ya descartada por ser par. \\
Luego, $5 \leq d$. \\
Por tanto, $d$ es de la forma $6k \pm 1$ para algún entero $k \geq 1$ y $5 \leq d \leq \sqrt{n}$.

\end{proof}

Entonces, dado $n$, el algoritmo que se desprende comprueba si 2 o 3 son divisores como proceso previo, si ello no ocurre, recorre los posibles divisores de la forma $6k \pm 1$ con $k \geq 1$. \\
Es decir, recorre -como posibles divisores- 5,7, \,11,13, \,17,19, etc., que son números de la forma $6k - 1$,$6k + 1$, con saltos de 6 en 6.  El Algoritmo~\ref{alg:Alg6k} implementa esta idea.

\begin{lstlisting}[caption={Testeo de primalidad optimizado mediante la regla del $6k \pm 1$}, label={alg:Alg6k}, mathescape=true]
# Entrada: n
# Salida: "True" cuando n es primo, "False" en caso contrario

n = 15                    # solo es un ej. puede cambiarse
primo_sw = True 
if (n $\mod$ 2) == 0 OR (n $\mod$ 3) == 0:
    primo_sw = False
else:
    for j = 5 hasta int($\sqrt{n}$) step 6:
        if (n $\mod$ j) == 0 OR (n $\mod$ (j+2)) == 0:
            primo_sw = False
            break              
print(primo_sw)
\end{lstlisting}

\section{Algoritmos de criba}
Son métodos deterministas basados en la exclusión iterativa y sistemática de múltiplos de números dentro de un  intervalo finito de los números naturales. En la criba de Eratóstenes se eliminan múltiplos de números primos, en variantes como la criba de Gries–Misra se consideran múltiplos tanto de primos como de compuestos. \\
Su fundamento descansa en la estructura del orden numérico y las propiedades de la divisibilidad: más que verificar la primalidad de un número aislado (aunque pueden usarse para ello), las cribas identifican y producen en bloque el conjunto de primos comprendidos entre 2 y $N$.

\subsection{Criba de Eratóstenes}
\label{subsec:Eratostenes}
Idea básica, dado un entero positivo $N$: \\ 
1. Se consideran todos los enteros 2, 3, 4, \dots, $N$. \\
2. Se marca el primer número primo 2 y se eliminan todos sus múltiplos estrictamente mayores que él, hasta $N$. \\  
3. Se pasa al siguiente entero no eliminado (que necesariamente es primo) y se eliminan todos sus múltiplos, como en el segundo paso. \\
Note que la criba usa un enfoque global y constructivo, eliminando de manera sistemática los múltiplos de cada primo, ya que todo $k$ compuesto, tiene un divisor primo $< k$ (incluso $< \sqrt{k}$).
\\  
4. El proceso continúa hasta que el primo actual $p$ satisface $p \leq \sqrt{N}$. \\  
5. Los enteros que no han sido eliminados al final del proceso son exactamente los números primos menores o iguales que $N$. 

\begin{ejemplo} Lista inicial con $N = 27$:\\
2, 3, 4, 5, 6, 7, 8, 9, 10, 11, 12, 13, 14, 15, 16, 17, 18, 19, 20, 21, 22, 23, 24, 25, 26, 27 \\
Primo actual $p = 2$. Tachamos sus múltiplos mayores:\\
2, 3, \cancel{4}, 5, \cancel{6}, 7, \cancel{8}, 9, \cancel{10}, 11, \cancel{12}, 13, \cancel{14}, 15, \cancel{16}, 17, \cancel{18}, 19, \cancel{20}, 21, \cancel{22}, 23, \cancel{24}, 25, \cancel{26}, 27 \\
Siguiente primo (actual) $p = 3$. Tachamos sus múltiplos mayores:\\
2, 3, \cancel{4}, 5, \cancel{6}, 7, \cancel{8}, \cancel{9}, \cancel{10}, 11, \cancel{12}, 13, \cancel{14}, \cancel{15}, \cancel{16}, 17, \cancel{18}, 19, \cancel{20}, \cancel{21}, \cancel{22}, 23, \cancel{24}, 25, \cancel{26}, \cancel{27}
Siguiente primo $p = 5$. Tachamos sus múltiplos mayores:\\
2, 3, \cancel{4}, 5, \cancel{6}, 7, \cancel{8}, \cancel{9}, \cancel{10}, 11, \cancel{12}, 13, \cancel{14}, \cancel{15}, \cancel{16}, 17, \cancel{18}, 19, \cancel{20}, \cancel{21}, \cancel{22}, 23, \cancel{24}, \cancel{25}, \cancel{26}, \cancel{27}
Siguiente primo $p = 7, 7 \nleq \sqrt{27}$. Fin. \\
Primos resultantes (números sin tachar) entre 2 y 27: 2, 3, 5, 7, 11, 13, 17, 19, 23.
\end{ejemplo}

Se eliminan los múltiplos de los primos $p$ tales que $p \leq \sqrt{N}$.\\
Todo los múltiplos del rango serán tachados si trabajamos sólo con esos primos. El siguiente Enunciado~\ref{enun:factorizacionAcotadaPorRaiz} justifica ello.

\begin{enunciado} [Factorización acotada por raíz]\label{enun:factorizacionAcotadaPorRaiz}\leavevmode\par Si $k$ es compuesto y $k \leq N$ entonces $k$ es múltiplo de un primo $p\leq \sqrt{N}$.
\end{enunciado}

\begin{proof}
El Enunciado~\ref{enun:cotaRaizdeN} muestra que un número compuesto $k$ tiene un divisor $d$ ($d \mid k$) tal que $1 < d \leq \sqrt{k}$. \\
Como $k \leq N$, es obvio que $\sqrt{k} \leq \sqrt{N}$.
Luego, $1 < d \leq \sqrt{N}$. \\
Si $d$ es primo entonces $p = d$ es un primo tal que $p \mid k$, es decir, $k$ es múltiplo de $p$ con $p \leq \sqrt{N}$. \\
Si $d$ no es primo, por el teorema fundamental de la aritmética, Enunciado~\ref{enun:teorFundAritm}, $d$ tiene una  factorización prima $d = p_1^{a_1} p_2^{a_2} \cdots p_k^{a_k}$. \\
Sea $p = p_1$. Es evidente que $p \mid d$. \\
$p \mid d$ y $d \mid k$, \setcounter{enumi}{5} entonces por el Enunciado~\ref{enun:propDivisibilidad}.\roman{enumi}: $p \mid k$. \\
Como $p \mid d$, \setcounter{enumi}{2} por el Enunciado~\ref{enun:propDivisibilidad}.\roman{enumi}: $p \leq d$. \\
$p \leq d$ y $d \leq \sqrt{N}$, luego $p \leq \sqrt{N}$. \\
Por lo tanto, $p$ es un primo tal que $p \mid k$, es decir, $k$ es múltiplo de $p$ con $p \leq \sqrt{N}$.
\end{proof}

Ello justifica detener el proceso una vez alcanzado el menor primo $p \nleq \sqrt{N}$. \\
En el ejemplo, nótese que cuando se trabaja con el primo $p = 3$, no hay necesidad de tachar el 6 pues éste ya ha sido tachado como múltiplo de 2. \\
En efecto, cuando consideramos el primo $p$ empezamos a tachar su múltiplos desde $p^{2}$ como justifica el siguiente Enunciado~\ref{enun:elimPrevia} cuya formalización coloquialmente dice: todos los múltiplos de $p$ que son menores que $p^{2}$ ya aparecen también como múltiplos de otros primos más pequeños.

\begin{enunciado} [Principio de eliminación previa]\label{enun:elimPrevia}\leavevmode\par Sea $p$ un número primo. $\forall m$ tal que $1 < m < p$, $\exists p'$ tal que $p'$ es primo, $p' \mid pm$ y $p' < p$.
\end{enunciado}

\begin{proof} Como $1 < m < p$, el número $m$ tiene un divisor primo $p' \leq m < p$. En efecto: \\
Por el teorema fundamental de la aritmética, Enunciado~\ref{enun:teorFundAritm}, $m$ tiene una  factorización prima $m = p_1^{a_1} p_2^{a_2} \cdots p_k^{a_k}$. \\
Sea $p' = p_1$. Es evidente que $p' \mid m$. \setcounter{enumi}{2} Por el Enunciado~\ref{enun:propDivisibilidad}.\roman{enumi}: $p' \leq m$. Como $m < p$ entonces $p' < p$. \\
Como $p' \mid m$, entonces \setcounter{enumi}{3}
por el Enunciado~\ref{enun:propDivisibilidad}.\roman{enumi}: $p' \mid pm$. \\
Así, $p'$ es el primo tal que $p' \mid pm$ y $p' < p$.
\end{proof}

\noindent El Algoritmo~\ref{alg:AlgCribaErat} implementa esta idea.
Note que es necesario utilizar una estructura (lista) de N posiciones que contenga True/False en la posición i si el i-ésimo número es o no primo.

\begin{lstlisting}[caption={Criba de Eratóstenes}, label={alg:AlgCribaErat}, mathescape=true]
# Entrada: N
# Salida: Lista de primos entre 0 y N 

N = 15                        # solo es un ej. puede cambiarse
Lista_primos = [True,$\dots$,True]   # lista de N elementos True
Lista_primos[0] = Lista_primos[1] = False  
p = 2
while p $\leq \lfloor \sqrt N \rfloor$:
    if Lista_primos[p]:
        # Tachamos todos los multiplos de p desde $p^{2}$ = p*p
        for i = $p^2$ hasta $N$ step p:
            Lista_primos[i] = False
    p = p+1
for i = 2 hasta $N$:
    if Lista_primos[i] == True:
        print(i)
\end{lstlisting}

\subsection{Criba de Eratóstenes con arreglo de bits y sólo impares}
Para mejorar la eficiencia en tiempo y memoria de la criba de Eratóstenes, se utiliza un arreglo de bits de tamaño $\lceil N/2 \rceil$ que representa exclusivamente a los números impares. Al ignorar los pares (excepto el 2, que se recibe un trato excepcional), reducimos el espacio a la mitad.

\begin{ejemplo} Se anotan los números impares y debajo el índice que tendrían en el arreglo de impares. \\
$
\begin{array}{ccccccccccccc}
3 & 5 & 7 & 9 & 11 & 13 & 15 & 17 & 19 & 21 & 23 & 25 & 27 \\
{\scriptscriptstyle 1} & {\scriptscriptstyle 2} & {\scriptscriptstyle 3} & {\scriptscriptstyle 4} & {\scriptscriptstyle 5} & {\scriptscriptstyle 6} & {\scriptscriptstyle 7} & {\scriptscriptstyle 8} & {\scriptscriptstyle 9} & {\scriptscriptstyle 10} & {\scriptscriptstyle 11} & {\scriptscriptstyle 12} & {\scriptscriptstyle 13}
\end{array}
$ \\
El ejemplo es sólo con fines de comprensión pues en vez de números impares las posiciones del arreglo tendrán en sus celdas bits que los representen y que serán interpretados como True (1) si el impar es primo o False (0) en otro caso.
\end{ejemplo} 

Existe una relación biunívoca entre el índice $i$ del arreglo y el número impar $m$: \\
De índice a número: $m = 2i + 1$. Por ejemplo, para $i = 4$, el número representado es $9$. \\
De número a índice: $i = \lfloor m/2 \rfloor$, pues un número impar m puede escribirse como 2i + 1, al dividirlo entre 2 y tomar la parte entera, se obtiene el índice i. Por ejemplo, para $m = 9$, calculamos su índice $i = \lfloor 9/2 \rfloor = 4$. \\
Cuando identificamos que un índice $i$ corresponde a un primo $p$ (que es impar por lo que estamos considerando), procedemos a tachar sus múltiplos en el arreglo. \\

Punto de partida: El primer múltiplo útil es $p^2$, que es impar por el Enunciado~\ref{enun:p2impar} y cuyo índice en el arreglo es $i = \lfloor p^2/2 \rfloor$. \\
Incremento: En la criba estándar se salta de $p$ en $p$. Dado que $p$ y $p^2$ son impares, entonces por el Enunciado~\ref{enun:parimpar} $p^2 + p$ es par, $p^2+p+p= p^2+2p$ es impar, $p^2 + 3p$ es par, etc.\\
Como el arreglo solo contiene impares, los múltiplos pares ($p^2 + p, p^2 + 3p$, etc.) no existen en la estructura. Por lo que  para saltar solo sobre los múltiplos impares ($p^2 + 2p, p^2 + 4p$, etc.), el incremento en números concretos es de $2p$, pero en el arreglo equivale a un salto de sólo $p$ pues sólo contiene impares. \\
Esta técnica permite mapear el espacio reducido del arreglo con los números concretos de forma exacta, logrando una criba más veloz y compacta. \\

El Algoritmo~\ref{alg:AlgCribaEratimparYbit} implementa esta idea.

\begin{lstlisting}[caption={Criba de Eratóstenes con sólo impares y bits}, label={alg:AlgCribaEratimparYbit}, mathescape=true]
# Entrada: N
# Salida: Lista de primos entre 0 y N 

N = 15                        # solo es un ej. puede cambiarse 
longitud_arreglo = $\lceil N/2 \rceil$ + ($N \mod 2$)
Lista_primos = [True,$\dots$,True]   # bitarray de $\lceil N/2 \rceil$ elementos True
Lista_Primos[0] = False # indice 0 corresponde al 1, que no es primo
limite = $\lfloor \sqrt N \rfloor$
for i=1 hasta $\lfloor limite/2 \rfloor$:
       if Lista_Primos[i] = True:
           p = 2$\cdot$i + 1
           inicio = $\lfloor p^2 / 2 \rfloor$
           for j=inicio hasta longitud_arreglo step p:
               Lista_Primos[j] = False
print(2)  # caso excepcional 2, unico primo par
for i=1 hasta longitud_arreglo: # longitud_arreglo-1 en index-0
       if Lista_Primos[i] = True:
           print(2$\cdot$i + 1)
\end{lstlisting}

\subsection{Criba incremental impar}
A diferencia de la criba clásica que tiene un arreglo booleano cubriendo todo el rango de 0 a N para el tachado masivo, aquí se tiene dos estructuras: \\
Una que contiene los primos hallados hasta el momento y que empieza con  2, el único primo par; en adelante sólo se incluirán números primos impares.
Otra que contiene los (siguientes) múltiplos pendientes de cada primo. No todos, uno a la vez.
Cuando $i=p$ y se detecta que $p$ es un número primo, se lo añade al conjunto de primos y se inicia su cadena de tachado añadiendo su primer múltiplo útil impar $p^2$ , este múltiplo es un número que recién será tratado cuando $i$ crezca hasta su valor. Además se incluye el factor primo del cual es múltiplo. \\
La lógica de tachado es distinta, es progresiva, paso a paso. Una vez que se detecta que $i=m$ es un múltiplo, se lo elimina de los múltiplos pendientes y se añade (actualiza) el siguiente múltiplo impar para cada factor primo del múltiplo eliminado.

Como sólo se añaden números primos impares entonces ignoraremos todos los números pares, ello significa que, como $m$ es impar y es múltiplo de $p$ con $p$ impar, el siguiente múltiplo de $p$ mayor a m es  $m+p$ que es par de modo que lo ignoramos, en cambio tomaremos a $m+2p$ como siguiente múltiplo (impar) de p mayor a $m$.

\begin{ejemplo} Supongamos dado un $N$ suficientemente grande. Al empezar: \\
Lista de primos: $[2]$. \\
Múltiplos pendientes (o compuestos):
\{ \} (la estructura está vacía).\\
$i = 3$. \\
No está entre los múltiplos pendientes, luego es primo, se añade a los primos. \\
Lista de primos: $[2, 3]$ \\
Se añade su primer múltiplo útil (impar), $3^2 = 9$, a los múltiplos pendientes, lo anotamos en el formato 9: [3]. \\
Múltiplos pendientes: \{ 9: [3] \} \\
Con $i = 5$ e $i = 7$ sucede lo mismo, no están entre los múltiplos pendientes, luego son primos, se añaden a los primos y se añade su cuadrado como primer múltiplo útil. \\
Lista de primos: $[2, 3, 5, 7]$ \\
Múltiplos pendientes: \{ 9: [3], 25: [5], 49: [7] \} \\
Siguiente número: $i = 9$. \\
Sí está entre los múltiplos pendientes con factor primo $3$. \\
Tachamos 9 (lo eliminamos de la estructura), y añadimos el siguiente múltiplo impar de $3$ mayor a 9: $9+2\cdot3=15$. \\
Múltiplos pendientes: \{ 15: [3], 25: [5], 49: [7] \} \\
$\dots$ \\
Supongamos que tenemos trabajado hasta el número 39, esta es la figura: \\
i = 39 \\     
Lista de primos: $[2, 3, 5, 7, 11, 13, 17, 19, 23, 29, 31, 37]$ \\
Múltiplos pendientes:
\{ 45: [5, 3], 49: [7], 121: [11], 169: [13], 289: [17], 361: [19], 529: [23], 841: [29], 961: [31], 1369: [37] \} \\
Siguiente $i$ (sólo impares): $i = 41$. \\
No está entre los múltiplos pendientes, luego es primo, se añade a los primos. \\
Lista de primos: $[2, 3, 5, 7, 11, 13, 17, 19, 23, 29, 31, 37, 41]$ \\
Se añade su primer múltiplo útil (impar), $41^2 = 1681$, a los múltiplos pendientes, en el formato 1681: [41], suponiendo que 1681 está dentro del rango N (en otro caso no se añade). \\
Múltiplos pendientes: \{ 45: [3, 5], 49: [7], 121: [11], 169: [13], 289: [17], 361: [19], 529: [23], 841: [29], 961: [31], 1369: [37], 1681: [41] \} \\
Siguiente número: $i = 43$. \\
No está entre los múltiplos pendientes, luego es primo, se añade a los primos: \\
Lista de primos: $[2, 3, 5, 7, 11, 13, 17, 19, 23, 29, 31, 37, 41, 43]$ \\
Y se añade $43^2 = 1849$ a los múltiplos pendientes (si está en el rango N):\\
Múltiplos pendientes: \{ 45: [3, 5], 49: [7], 121: [11], 169: [13], 289: [17], 361: [19], 529: [23], 841: [29], 961: [31], 1369: [37], 1681: [41], 1849: [43] \} \\
Siguiente número: $i = 45$. \\
Sí está entre los múltiplos pendientes con factores primos $3$ y $5$, en el formato 45: [3, 5]. \\
Tachamos 45 (lo eliminamos de la estructura), y añadimos los siguientes múltiplos impares de $3$ y $5$, suponiendo que está(n) dentro del rango N (en otro caso no se añade(n)): $45+2\cdot3=51$ para 3 y $45+2\cdot5=55$ para 5. Este es el resultado: \\
Múltiplos pendientes: \{ 49: [7], 51: [3], 55: [5], 121: [11], 169: [13], 289: [17], 361: [19], 529: [23], 841: [29], 961: [31], 1369: [37], 1681: [41], 1849: [43] \} \\
Los primos quedan tal cual estaban: \\
Lista de primos: $[2, 3, 5, 7, 11, 13, 17, 19, 23, 29, 31, 37, 41, 43]$.
\end{ejemplo}

Es necesarios guardar todos los factores primos de un múltiplo pendiente para mantener la integridad de la criba a medida que crece. \\
En efecto, la criba incremental no es solo una lista de números, sino un sistema dinámico donde cada número primo $p$ genera su propia progresión aritmética de múltiplos impares:  $\{p^2, p^2 + 2p, p^2 + 4p, \dots \}$. \\
Un número compuesto $m$ es, por definición, el punto de intersección de dichas progresiones de sus factores primos. \\
Si sólo almacenáramos un factor (por ejemplo el más pequeño) en cada múltiplo las progresiones  de los demás factores se detendrían. Al no actualizarse las progresiones de los factores omitidos, no se registran sus múltiplos futuros, clasificando erróneamente números compuestos como primos.

\begin{ejemplo}
$i =63$. Deberíamos tener 63: [3, 7] en los múltiplos pendientes. \\
Tachamos $63$ y añadimos $ 69 = 63+2\cdot 3$ y $77 = 63+2\cdot 7$ como siguientes múltiplos de $3$ y $7$. \\
Supongamos que sólo registramos el factor primo $3$ para $63$, en el formato 63: [3]. \\
La progresión del 7 se verá detenida desde el 63. \\
Cuando lleguemos a $i=77$ múltiplo de 7, no habrá ningún registro esperándolo, no estárá entre los múltiplos pendientes, conluyendo erróneamente que el 77 es primo. \\
\end{ejemplo}

El Algoritmo~\ref{alg:AlgCribaIncremental} implementa esta idea.

\begin{lstlisting}[caption={Criba incremental con impares}, label={alg:AlgCribaIncremental}, mathescape=true]
# Entrada: N
# Salida: Lista de primos entre 0 y N 

N = 15                        # solo es un ej. puede cambiarse 
Lista_primos = [2]            # el unico par (trato excepcional)
compuestos = {}
for i=3 hasta N step 2:       # solo impares
    if i not in compuestos:
        agregar i a Lista_primos
        if $i^2 \leq$ N:                # solo si esta en el rango
            agregar $i^2$:[i] a compuestos
    else:
        for p in compuestos[i]:
            # como p es impar, avanzamos de 2p en 2p
            sgte = i+2$\cdot$p
            if sgte $\leq$ N: 
                if sgte in compuestos:
                    agregar p a la lista asociada a sgte
                else: 
                    agregar sgte:[p] a compuestos
        tachar compuestos[i]   # eliminar de compuestos
print(Lista_primos)
\end{lstlisting}

Tener dos estructuras que crecen progresivamente, eliminando sus elementos iniciales poco a poco en los múltiplos pendientes puesto que solo se mantienen los múltiplos activos, evita tener estructuras masivas desde el inicio. \\
También evita tener que recorrer listas largas de múltiplos desde el comienzo y permite liberar memoria de manera progresiva conforme los compuestos dejan de ser relevantes. \\
Además, si guardamos los últimos estados de ambas estructuras, el algoritmo puede extenderse a valores mayores de N sin necesidad de rehacer todo el cribado, ya que la propagación se mantiene local y dinámica. Es lo que Bengelloun \cite{ DBLP:journals/acta/Bengelloun86} llama generación infinita de primos.

\subsection{Criba por rueda (wheel sieve)}
La idea es evitar la consideración de números que ya sabemos que son múltiplos de un grupo dado de primos.\\
Sea $S = \{p_1, p_2, \dots, p_k\}$ el conjunto de los primeros $k$ números primos.\\

Sea que estamos trabajando con los números entre 2 y N. \\
En vez de tachar, como en la criba clásica, los múltiplos de $p_1=2$, de $p_2=3$, etc., evitaremos mucho trabajo obviando a sus múltiplos de una sola vez de un modo muy ingenioso.

\begin{ejemplo}
Con $S = \{2,3,5\}$ y $N = 27$:\\
2, 3, 4, 5, 6, 7, 8, 9, 10, 11, 12, 13, 14, 15, 16, 17, 18, 19, 20, 21, 22, 23, 24, 25, 26, 27 \\
Evitaremos considerar números como 4, 9 o 25 que se  sabe que son múltiplos de algún elemento de S.
\end{ejemplo}

Al producto $M = p_1p_2 \dots p_k$ lo denominaremos primorial de orden $k$.\\
Sea $CRR(M)$ el conjunto reducido de restos de $M$ (Definición~\ref{def:crr}). \\
Los factores primos de $M$ son los que están en S.

\begin{enunciado} [Principio de la rueda por residuos]\label{enun:ppioRuedaResid}\leavevmode\par Un número entero $x$ es coprimo con $M$ si y sólo si al dividirlo por $M$ su residuo pertenece al conjunto $CRR(M)$. Sea $x \bmod M = r$, el enunciado puede reescrbirse así: \\
$mcd(x,M) = 1$ si y sólo si $r \in CRR(M)$.
\end{enunciado}

\begin{proof}
Sea $x \bmod M = r$. \\
Por el algoritmo de la división: \\
$x = qM + r$, donde $0 \le r < M$. Todo número x tiene esta forma.\\
\setcounter{enumi}{7} 
El Enunciado~\ref{enun:propMcd}.\roman{enumi} establece que $mcd(x,M) = mcd(r,M)$.\\
Si $mcd(x, M) = 1$, entonces $mcd(r,M) = 1$, lo que implica que $r \in CRR(M)$ por definición de $CRR(M)$. \\
Si $r \in CRR(M)$, entonces $mcd(r,M) = 1$, por lo tanto $mcd(x, M) = 1$.
\end{proof}

El Enunciado~\ref{enun:ppioRuedaResid} permite una idea ingeniosa. \\

Primera parte de la ingeniosa idea: \\
Cualquier número $x$ que dé un residuo $r \in CRR(M)$ al dividirlo entre $M$ es coprimo con $M$, es decir, $mcd(x,M)=1$ y puede escribirse como:  $k'M + r_i$, con $r_i \in CRR(M)$. \\
Esto significa que $x$ no es divisible por ningún primo de $S$, puesto que si algún $p_i \in S$ fuese tal que $p_i \mid x$, como $p_i \mid M$, resulta que $mcd(x,M) > 1$. Pero éste no es el caso. \\
Al no ser x divisible por ningún primo de $S$, se convierte en un candidato a primo nuevo (nuevo respecto de los que están en $S$).

\begin{ejemplo}
Sea $S = \{2,3,5\}$. Sea $M=2 \cdot 3 \cdot 5 = 30$.\\
$CRR(M)=\{1, 7, 11, 13, 17, 19, 23, 29\}$, cualquiera de estos números es coprimo con $M$.\\
Con $x=49=1 \cdot 30 + 19$, $mcd(49,30)=mcd(19,30)=1$.\\
Ni 2, ni 3, ni 5 dividen a 49, es un candidato a primo nuevo (respecto de S). \\
Con $x=83=2 \cdot 30 + 23$, $mcd(83,30)=mcd(23,30)=1$.\\
Ni 2, ni 3, ni 5 dividen a 83, es un candidato a primo nuevo (respecto de S).
\end{ejemplo}

Segunda parte de la ingeniosa idea: \\
Cualquier número $x$ que dé un residuo $r$ que no esté en $CRR(M)$, es decir, que no sea coprimo con $M$, es decir, tal que $mcd(x,M)=mcd(r,M)>1$, al compartir $r$ un factor primo con $M$, entonces $x$ también lo comparte. Luego, $x$ es divisible por algún primo de $S$ y no puede ser un primo nuevo (respecto de $S$). \\
Dicho de otro modo, ningún número no coprimo con $M$ puede escribirse como  $k'M + r_i$, con $r_i \in CRR(M)$.\\
Excluiremos estos números de nuestras consideraciones.

\begin{ejemplo}
Sea $S = \{2,3,5\}$. Sea $M=2 \cdot 3 \cdot 5 = 30$.\\
Con $x=28$, $mcd(28,30)=2 > 1$. Por tanto $2 \mid x$. \\
Con $x=65=2 \cdot 30 + 5$, $mcd(65,30)=mcd(5,30)=5 > 1$. Por tanto $5 \mid x$. \\
Con $x=90=3 \cdot 30 + 0$, $mcd(90,30)=mcd(0,30)=30 > 1$. Como $2 \mid 30$ y $30 \mid x$: $2 \mid x$. \\
Con $x=363=12 \cdot 30 + 3$, $mcd(363,30)=mcd(3,30)=3 > 1$. Por tanto $3 \mid x$. 
\end{ejemplo}

Respecto del $CCR_{canonico}(M)$, lo que se se dijo en la primera y segunda parte de la idea ingeniosa, es un patrón que se repite cada $M$ números. \\
Por eso hablamos de una rueda: cada vuelta reproduce exactamente los mismos residuos: los que no nos interesan y los que sí.\\

Tercera parte de la ingeniosa idea: \\
Dado $M$. \\
En vez de tener todos los números de 2 a $N$ para añadir los que son primos a la lista de primos o tachar los compuestos, primero obtendremos el $CRR(M)$.\\
Luego, se generan números x del tipo: $k'M + r_i$, con $r_i \in CRR(M)$.  \\
Al variar $k'$, se generan todos los números coprimos con $M$ y de ahí que son buenos candidatos, pues no son múltiplos (al menos) de los primos de $S$.\\
Obviamente se respeta el límite $N$ y obviamente es mejor ignorar los candidatos que no sean mayores al más grande primo de $S$. \\
Se dice que se generan candidatos por clases de residuos pues dichos candidatos pertenecen a familias enteras de números que comparten la misma clase de residuo respecto de $M$.

\begin{ejemplo}
$N=23$. Rueda de orden $k=2$.\\
$S = \{2, 3\}$. $M = 2 \cdot 3 = 6$. $CRR(6) = \{1, 5\}$.\\
Generar candidatos de la forma $k' \cdot M + r_i$ con $r_i \in CRR(M)$: \\
Para $k'=0$: $\{1, 5\}$ (se ignora el $1$, no es mayor a $3$)\\
Para $k'=1$: $\{6+1, 6+5\} = \{7, 11\}$ \\
Para $k'=2$: $\{12+1, 12+5\} = \{13, 17\}$ \\
Para $k'=3$: $\{18+1, 18+5\} = \{19, 23\}$  Para $k' > 3$ excedemos el rango\\
candidatos = $\{5, 7, 11, 13, 17, 19, 23\}$ \\
La clase del residuo $1$ es $[1] = \{1,7,13,19, \cdots \}$ \\
La clase del residuo $5$ es $[5] = \{5,11,17,23, \cdots \}$ \\
Sólo anotamos y tomamos en cuenta los que están dentro del rango $N$.
\end{ejemplo}

Nótese que en el fondo estamos manejando bloques $[k'M,(k'+1)M[$ de tamaño M, es decir, con M enteros. \\
Cada bloque contiene exactamente los mismos residuos coprimos trasladados por múltiplos de M, su cantidad es $CRR(M)$. \\
Para cubrir el rango de $0$ a $N$ con estos bloques de tamaño $M$, el último bloque que puede aportar candidatos es aquel con índice
$k'=\left\lfloor \frac{N}{M}\right\rfloor$ porque cualquier $k'>\frac{N}{M}$ produciría $k'M+r_i>N$, fuera del rango. \\

En términos prácticos, la rueda es una optimización: reduce el espacio de búsqueda y el número de operaciones ya que evita considerar números descartados de antemano -que sabemos compuestos- y se concentra únicamente en los que tienen posibilidades de ser primo pues se aplica sobre un conjunto compacto de candidatos que se generan en orden creciente. \\
Es claro que puede ser necesario un manejo adecuado de listas o estructuras donde se almacenen: \\
- Lista de primos. \\
- Lista candidatos. \\
- Los candidatos que resulten tachados al ser compuestos. Lista de tachados. \\
En lugar de revisar todos los números (como en la criba clásica), solo trabajamos con los que están entre los candidatos, que son aquellos cuyo residuo módulo $M$ está en $CRR(M)$, es decir, son coprimos con M. \\
Ello elimina de una sola vez todos los múltiplos de S.\\
Con el candidatos $num$ se procede así: \\
- Si no está en la lista de tachados, se añade al conjunto de primos, luego se tachan sus múltiplos a partir de $num^2$ y dentro del límite $N$, incrementando de $num$ en $num$, siempre que dichos múltiplos además estén en la lista de candidatos. Tachar un múltiplo que está entre los candidatos es equivalente a añadirlo a la lista de tachados (puede anularse de la lista de candidatos pero, como se va en orden creciente, no es necesario salvo para liberar memoria). \\
- Si está en tachados, se ignora pues es compuesto. \\
Se procede así con cada elemento de la lista de candidatos.

\begin{ejemplo}
$N=26$. Rueda de orden $k=2$.\\
$S = \{2, 3\}$. $M = 2 \cdot 3 = 6$. $CRR(6) = \{1, 5\}$.\\
candidatos = $\{5, 7, 11, 13, 17, 19, 23, 25\}$ \\
ListaPrimos = [ ] \\
tachados = \{ \} \\
$num = 5$ \\
No está en tachados, se agrega a la lista de primos. \\
Tachamos su múltiplos desde $5^2 = 25$. Agregamos $25$ a la lista de tachados pues $25$ también está entre los candidatos. Como $25+5 = 30 > N$, nos detenemos. \\
ListaPrimos = [5] \\
tachados = \{25\} \\
$num = 7$ \\
No está en tachados, se agrega a la lista de primos. \\
Tachamos su múltiplos desde $7^2 = 49$. No sólo $49$ no está en candidatos sino que $49 > N$, nos detenemos. \\
$ListaPrimos = [5,7]$ \\
tachados = \{25\} \\
Con $num = p = 11, 13, 17, 19, 23$ sucede los mismo: \\
No están en tachados, se añaden a la lista de primos. En cada caso su primer múltiplo útil, $p^2$, no sólo no está entre los candidatos sino que queda fuera del rango $N$, no se agrega nada a tachados.\\
$ListaPrimos = [5,7, 11, 13, 17, 19, 23]$ \\
tachados = \{25\} \\
$num = 25$ sí está en tachados, se ignora. \\
Hemos acabado con todos los candidatos. \\
Como salida, la lista de primos se une a los primos base S. \\
Nótese que este ejemplo es equivalente al algoritmo que hemos denominado por división tentativa $6k \pm 1$.
\end{ejemplo}

Todas estas ideas se implementan en el Algoritmo~\ref{alg:AlgCribaRueda}.

\begin{lstlisting}[caption={Criba por rueda}, label={alg:AlgCribaRueda}, mathescape=true]
# Entrada: N, k
# Salida: Lista de primos entre 0 y N

N = 15                        # solo es un ej. puede cambiarse
Lista_primos = []
Lista_tachados = []
k = 3                         # orden  de la rueda
$M = p_1 \cdot p_2 \cdot \, \cdots \, \cdot p_k$
$S = primosbase = [p_1, p_2, \cdots , p_k]$
$residuos = \{ r \; / \; 1 \leq r < M, mcd(r,M) = 1 \}$   # CRR(M)

Lista_primos = primosbase
candidatos   = $[k' \cdot M + r : k' \geq 0, \, r \in residuos, max(S) < k' \cdot M+r \leq N]$

for num in candidatos:           # en orden creciente
    if num not in Lista_tachados:
        agregar num a Lista_primos
        for m desde $num^{2}$ hasta N step $num$:
            if $m \in candidatos$:
                agregar m a Lista_tachados
print(Lista_primos)
\end{lstlisting}

La estructura de la rueda es una enumeración por clases de residuos: $k'M + r_i$, con $r_i \in CRR(M)$. \\
Nótese que los múltiplos de un nuevo primo p ($num$ en el algoritmo) se tachan si están en la lista de candidatos que es generada según la filosofía de la rueda, pero dichos múltiplos se generan de un modo ajeno a dicha filosofía y más como en una criba clásica, desde su cuadrado, de p en p.

\subsection{Criba de Gries y Misra}
\label{sec:Gries}
Note que para una lista inicial con $N = 15$:\\
2, 3, 4, 5, 6, 7, 8, 9, 10, 11, 12, 13, 14, 15 \\
En la criba de Eratóstenes, el 12 se tacha cuando se trabaja con los primos 2 y 3. Redundante.

Gries y Misra \cite{Gries1978LinearSA} eliminan la redundancia mencionada al precio de tener estructuras de datos adicionales. Lo interesante es que los autores hallan la forma de que cada número compuesto sea eliminado de la lista por uno y solo uno de sus divisores, su mínimo factor primo. En el caso de la lista anterior, el 12 se tacha cuando se trabaja con $i = 6$ y $p = 2$, no se tacha más ni antes ni después. \\ 
Eso hace que esta criba tenga un comportamiento lineal. \\

Para cada $N$ se trabaja con una lista del menor factor primo (mfp) para cada número entre 2 y $N$, la anotaremos entre paréntesis. \\ 
Se trabaja también con una lista de los primos hallados hasta el momento. \\ 
Por comodidad de lectura en los ejemplos añadiremos una lista de los números entre 2 y $N$ entre llaves “\{“ y “\}”, aunque no es en absoluto necesario. \\ 
Con miras a la implementación incluimos las posiciones 0 y 1 en cada lista, pero las ignoramos con un guión bajo “\_”. \\

\begin{ejemplo}
Para $N = 15$, inicialmente la lista del menor factor primo tiene sólo ceros: \\ 
\{\_, \_, 2, 3, 4, 5, 6, 7, 8, 9, 10, 11, 12, 13, 14, 15\} \\ 
(\_, \_, 0, 0, 0, 0, 0, 0, 0, 0, 00, 00, 00, 00, 00, 00) 
\end{ejemplo}

La idea es esta:  \\ 
Para cada $i \in $ \{2, 3, 4, …, $N$ \}: \\ 
Si su menor factor primo está en 0, lo reemplazamos por $i$ y lo agregamos a la lista de primos. \\ 
Luego, para todos los primos $p_j$ que no sean mayores al menor factor primo de $i$, hallamos $(p_j \cdot i)$: si $(p_j \cdot i) \leq N$ y el menor factor primo de $(p_j \cdot i)$ es cero, lo reemplazamos por $p_j$. \\
(lo que hacemos es poner como menor factor primo del número compuesto $(p_j \cdot i)$ a su divisor $p_j$; note que si no es cero, ya pusimos un divisor antes y ya no se hace nada, poner un divisor es el equivalente a tachar). \\

\begin{ejemplo} Con $N = 15$. \\
i=2 \\ 
\{\_, \_, 2, 3, 4, 5, 6, 7, 8, 9, 10, 11, 12, 13, 14, 15\} \\ 
(\_, \_, 2, 0, 2, 0, 0, 0, 0, 0, 00, 00, 00, 00, 00, 00) \\ 
primos = \{2\} \\ 
i=3 \\ 
\{\_, \_, 2, 3, 4, 5, 6, 7, 8, 9, 10, 11, 12, 13, 14, 15\} \\ 
(\_, \_, 2, 3, 2, 0, 2, 0, 0, 3, 00, 00, 00, 00, 00, 00) \\ 
primos = \{2,3\} \\ 
Expliquemos un poco este paso. \\
En el paso anterior añadimos 2 como factor primo menor de los enteros 2 y 4. \\
Estamos trabajando con $i = 3$, su factor primo menor $p$ está en cero $p = 0$, por lo tanto lo reemplazamos por $p =3$. Ahora tenemos esta lista de primos: \{2,3\}. \\
Como para $p_j = 2$ y $p_j = 3$: $p_j \leq p$, obtenemos $(p_j \cdot i)$, es decir, $(p_j \cdot 3)$ y en el número resultante ponemos $p_j$ como su factor primo menor. Así, $(2 \cdot 3) = 6$, actualizamos el factor primo menor de 6 a 2; $(3 \cdot 3) = 9$, actualizamos el factor primo menor de 9 a 3. \\
i=4 \\ 
\{\_, \_, 2, 3, 4, 5, 6, 7, 8, 9, 10, 11, 12, 13, 14, 15\} \\ 
(\_, \_, 2, 3, 2, 0, 2, 0, 2, 3, 00, 00, 00, 00, 00, 00) \\ 
primos = \{2,3\} \\ 
i=5 \\ 
\{\_, \_, 2, 3, 4, 5, 6, 7, 8, 9, 10, 11, 12, 13, 14, 15\} \\ 
(\_, \_, 2, 3, 2, 5, 2, 0, 2, 3,  02, 00, 00, 00, 00, 03) \\ 
primos = \{2,3,5\} \\ 
i=6 \\ 
\{\_, \_, 2, 3, 4, 5, 6, 7, 8, 9, 10, 11, 12, 13, 14, 15\} \\ 
(\_, \_, 2, 3, 2, 5, 2, 0, 2, 3,  02, 00, 02, 00, 00, 03) \\ 
primos = \{2,3,5\} \\ 
Expliquemos un poco este paso. \\
En algún paso anterior añadimos $p = 2$ como factor primo menor de 6, que no es cero, por lo tanto se queda así. \\
Estamos trabajando con $i = 6$ y esta lista de primos: \{2,3,5\}. \\
Como $p_j = 2$ es el único primo tal que $p_j \leq p$, obtenemos $(p_j \cdot 6)$, es decir, $(2 \cdot 6) = 12$ y en el número resultante 12 ponemos a 2 como su factor primo menor. \\
No trabajamos con $p_j = 3$ ni $p_j = 5$ pues en estos casos no se cumple que $p_j \leq p$. \\
i=7 \\ 
\{\_, \_, 2, 3, 4, 5, 6, 7, 8, 9, 10, 11, 12, 13, 14, 15\} \\ 
(\_, \_, 2, 3, 2, 5, 2, 7, 2, 3,  02, 00, 02, 00, 02, 03) \\ 
primos = \{2,3,5,7\} \\ 
i=8 \\ 
\{\_, \_, 2, 3, 4, 5, 6, 7, 8, 9, 10, 11, 12, 13, 14, 15\} \\ 
(\_, \_, 2, 3, 2, 5, 2, 7, 2, 3,  02, 00, 02, 00, 02, 03) \\ 
primos = \{2,3,5,7\} \\ 
i=9 \\ 
\{\_, \_, 2, 3, 4, 5, 6, 7, 8, 9, 10, 11, 12, 13, 14, 15\} \\ 
(\_, \_, 2, 3, 2, 5, 2, 7, 2, 3,  02, 00, 02, 00, 02, 03) \\ 
primos = \{2,3,5,7\} \\ 
i=10 \\ 
\{\_, \_, 2, 3, 4, 5, 6, 7, 8, 9, 10, 11, 12, 13, 14, 15\} \\ 
(\_, \_, 2, 3, 2, 5, 2, 7, 2, 3,  02, 00, 02, 00, 02, 03) \\ 
primos = \{2,3,5,7\} \\ 
i=11 \\ 
\{\_, \_, 2, 3, 4, 5, 6, 7, 8, 9, 10, 11, 12, 13, 14, 15\} \\ 
(\_, \_, 2, 3, 2, 5, 2, 7, 2, 3,  02, 11, 02, 00, 02, 03) \\ 
primos = \{2,3,5,7,11\} \\ 
i=12 \\ 
\{\_, \_, 2, 3, 4, 5, 6, 7, 8, 9, 10, 11, 12, 13, 14, 15\} \\ 
(\_, \_, 2, 3, 2, 5, 2, 7, 2, 3,  02, 11, 02, 00, 02, 03) \\ 
primos = \{2,3,5,7,11\} \\ 
i=13 \\ 
\{\_, \_, 2, 3, 4, 5, 6, 7, 8, 9, 10, 11, 12, 13, 14, 15\} \\ 
(\_, \_, 2, 3, 2, 5, 2, 7, 2, 3,  02, 11, 02, 13, 02, 03) \\ 
primos = \{2,3,5,7,11,13\} \\ 
i=14 \\ 
\{\_, \_, 2, 3, 4, 5, 6, 7, 8, 9, 10, 11, 12, 13, 14, 15\} \\ 
(\_, \_, 2, 3, 2, 5, 2, 7, 2, 3,  02, 11, 02, 13, 02, 03) \\ 
primos = \{2,3,5,7,11,13\} \\ 
i=15 \\ 
\{\_, \_, 2, 3, 4, 5, 6, 7, 8, 9, 10, 11, 12, 13, 14, 15\} \\ 
(\_, \_, 2, 3, 2, 5, 2, 7, 2, 3,  02, 11, 02, 13, 02, 03) \\ 
primos = \{2,3,5,7,11,13\}
\end{ejemplo}

Note que como cada compuesto tiene asociado su menor factor primo, su factorización se convierte en un rastreo lineal y determinista con divisiones predichas. \\
En el ejemplo, 14 es compuesto y su menor factor primo es 2, 14/2 = 7, 7 es primo pues su menor factor primo es él mismo. Luego 14 = 2 $\cdot$ 7.\\
\\
Es usual llamar bucle externo al recorrido de los número $i$ de 2 a $N$. \\
Es usual llamar bucle interno al recorrido de la lista de primos (que van en orden creciente). \\

Es necesario justificar el procedimiento.

\begin{enunciado}[Estabilidad del menor factor primo]\label{enun:estabilidadMenorPrimo}\leavevmode\par
Sea $i \in \Z$ \, un entero positivo cuyo factor primo menor es $p$. Sea $p_j$ un número primo tal que $p_j \leq p$. Entonces el factor primo menor de $m = p_j \cdot i$ es exactamente $p_j$.        
\end{enunciado}

\begin{proof} De $m = p_j \cdot i$ sabemos que $p_j \mid m$ y por hipótesis $p_j$ es primo. \\
Supongamos que existe un primo $q < p_j$ tal que $q \mid m$, es decir, $q \mid (p_j \cdot i)$. \\
\setcounter{enumi}{2}
Como $q$ y $p_j$ son primos es obvio que $q \nmid p_j$, entonces por el Enunciado~\ref{enun:propDivprim}.\roman{enumi}, $q \mid i$. \\
Si $q \mid i$ y $q$ es primo, entonces $q$ es un factor primo de $i$. \\
Por hipótesis, sabemos que $p$ es el menor factor primo de $i$, por lo tanto $q \geq p$. \\
Nuestro supuesto dice que $q < p_j$. Por hipótesis sabemos que  $p_j \leq p$, luego  $q < p$\\
Resumiendo: $q \geq p$ y $q < p$: una contradicción. \\
De ahí que nuestro supuesto es erróneo, lo que prueba el enunciado.
\end{proof}

Con ello se garantiza que la asignación del menor factor primo para cada número compuesto es correcta. \\

\begin{enunciado}[Caracterización única de compuestos en Gries–Misra]\label{enun:compuestoUnico}\leavevmode\par
Sea $ N \geq 3$. $\forall m$ tal que $m$ es compuesto y $2 < m \leq N$, existe un único par $(p,i)$ tal que: \\
a) $p$ es el menor factor primo de $m$. \\
b) $i \in \N$ y $m = p \cdot i$. \\
c) El algoritmo de Gries–Misra marca $m$ como compuesto exactamente cuando el bucle externo toma el valor $i$ y el bucle interno procesa el primo $p$.
\end{enunciado}

\begin{proof} Denotaremos al menor factor primo de $x$ como $\mathrm{mfp}(x)$. \\
Llamaremos marcado de $m$ al hecho de anotar su $\mathrm{mfp}(m)$ en la lista $\mathrm{mfp}$ de menores factores primos. \\
a) Sea $m > 2$ un entero compuesto. Por el teorema fundamental de la aritmética, $m$ admite una factorización prima única. Sea $p =$ mfp($m$). Definimos $i = m / p$.\\
Entonces $i \in \N$ e $i \geq 2$, ya que si $i = 1$ se tendría $m = p$, contradiciendo que $m$ es compuesto.
Además, por la minimalidad de $p$, todo factor primo de $i$ es mayor o igual que $p$; en particular, $p \leq \mathrm{mfp}(i)$. \\
La unicidad del par $(p,i)$ es inmediata: $p$ está determinado de manera única como el menor factor primo de $m$, e $i$ queda fijado como $m/p$. Recibe el nombre de par canónico de $m = (\mathrm{mfp}(m),m / \mathrm{mfp}(m)) = (p,i)$. \\ \\
b) El algoritmo de Gries–Misra recorre todos los valores de $i$ desde 2 hasta $N$. \\
Para cada valor de $i$, el bucle interno itera sobre los primos $p_j$ ya identificados, en orden creciente, mientras se cumpla la condición $p_j \leq \mathrm{mfp}(i)$. \\
Como hemos visto en el inciso anterior, para el valor $i = m / p$ se cumple $p \leq \mathrm{mfp}(i)$. Por tanto, cuando el bucle externo alcanza este valor de $i$, el bucle interno necesariamente procesará el primo $p$. \\
Además, $p$ ya ha sido identificado como primo antes o en ese momento, pues $p \leq i$ y el algoritmo descubre los primos en orden creciente. \\ \\
c) En la iteración en la que el bucle externo toma el valor $i$ y el bucle interno procesa el primo $p$, el algoritmo evalúa el producto $(p \cdot i) = m$ y marca $m$ como compuesto, colocando $p$ en la lista mfp como el menor factor primo de $m$.
\end{proof}

Con ello se garantiza que cada compuesto se marca exactamente una vez, en la iteración correspondiente, pues el inciso a) muestra la existencia y unicidad del par $(p,i)$, el inciso b) que dicho par se alcanza en el algoritmo y el inciso c) que $m$ se marca a través de ese par único, al menos con la restricción de trabajar con los primos $p_j$ tales que $p_j \leq p$, dado $i$ y $\mathrm{mfp}(i) = p$.

\begin{enunciado}[Redundancia del marcado no canónico en Gries–Misra]\label{enun:marcadoRedundante}\leavevmode\par
Sea $i \geq 2$  y sea $q$  un primo tal que $q > \mathrm{mfp}(i)$. Si el algoritmo evalúa el producto $m = i \cdot q$ entonces el marcado de $m$ como compuesto es redundante pues $m$ será marcado en la iteración correspondiente a su par canónico $(p,m/p)$.
\end{enunciado}

\begin{proof}
Sea $p = \text{mfp}(i)$ y sea $q > p$ primo. Definimos $m = i \cdot q$. \\
$\text{mfp}(m) = \min(q, \text{mfp}(i)) = \min(q, p) = p$. \\
El par canónico es $(p,i')$ donde $i' = \frac{m}{p} = \frac{i \cdot q}{p}$. \\
Al ser $p$ un factor primo de $i$, es claro que $p \mid i$, escribimos $i = p \cdot k$ con $k \geq 1$. Entonces:
$i' = \frac{p \cdot k \cdot q}{p} = k \cdot q$. \\
Comparando $i'$ con $i$, $i' = k \cdot q = \frac{i}{p} \cdot q = i \cdot \frac{q}{p} > i$ pues $q > p$ por hipótesis. \\
Por tanto, el bucle externo alcanza $i$ antes que $i'$. Esta es una de las razones por las que $m$ será marcado luego en la iteración $i'$, no en $i$.\\
Cuando el bucle externo alcance $i'$: \\
- Caso $k = 1$: $i' = q$, así que $\text{mfp}(i') = q > p$. Por tanto, se cumple $\text{mfp}(i') \geq p$. \\
- Caso $k > 1$: $\text{mfp}(i') = \min(\text{mfp}(k), q)$. \\
Si $\text{mfp}(k) > q$, entonces $\min(\text{mfp}(k), q) = q$ y como $q > p$, entonces $\min(\text{mfp}(k), q) > p$, y por tanto $\min(\text{mfp}(k), q) \geq p$. \\
Si $q \geq \text{mfp}(k)$, entonces $\min(\text{mfp}(k), q) = \text{mfp}(k)$. Pero $\text{mfp}(k) \geq p$ (por minimalidad de $p$ en $i = p \cdot k$), luego $\min(\text{mfp}(k), q) \geq p$. \\
En ambos casos, $\text{mfp}(i') \geq p$, garantizando que el bucle interno en $i'$ procesará el primo $p$ y marcará $m = p \cdot i'$ como compuesto. Es aquí donde $m$ se marcará canónicamente.\\
Así, cualquier marcado de $m$ mediante el par $(q, i)$ con $q > \text{mfp}(i)$ sería redundante, pues $m$ inevitablemente será marcado en la iteración canónica posterior $(p, i')$. La condición $p_j \leq \text{mfp}(i)$ del algoritmo es por tanto segura y óptima.
\end{proof}

Con ello se garantiza que todo compuesto se marque únicamente en su par canónico y en ninguna otra iteración.
Cualquier par $(q,i)$ con $q > \text{mfp}(i)$ no es canónico por lo que el marcado del compuesto $m = q \cdot i$ sería redundante y es bueno que quede excluido por la condición $p_j \leq \text{mfp}(i)$ pues no se trata solo de evitar duplicaciones, permitirlo -violando la condición del bucle interno- implicaría asignar a $m$ un primo que no es su menor factor primo, desvirtuando el invariante fundamental del algoritmo, según el cual $\text{mfp}(m)$ se asigna exactamente una vez y correctamente.\\

Todas estas ideas se implementan en el Algoritmo~\ref{alg:AlgCribaGries}.

\begin{lstlisting}[caption={Criba de Gries-Misra}, label={alg:AlgCribaGries}, mathescape=true]
# Entrada: N
# Salida: Lista de primos entre 0 y N
#         Lista del menor factor primo de m para $2 \leq m \leq N$ 

N = 15                        # solo es un ej. puede cambiarse
Lista_primos = []
mfp = [0,$\dots$,0]                  # lista de N elementos cero
for i = 2 hasta N:
    if mfp[i] == 0:
        mfp[i] = i
        agregar i a Lista_primos
    for p in Lista_primos:
        if (p > mfp[i]) OR (i $\cdot$ p) > n:
            break
        mfp[i $\cdot$ p] = p
print(Lista_primos)
print(mfp)
\end{lstlisting}

Cuando el menor factor primo de $i$ aparece como cero,  hay ausencia de una descomposición previa, no se pudo haber cambiado su menor factor primo cero en una iteración anterior, puesto que no posee divisores propios. En el recorrido, el procedimiento hallará infaliblemente a $i$ con su menor factor primo como cero, lo identifica como primo y por tanto es también su propio factor primo menor, cambiando cero por $p = i$. \\
Además se incorpora $i$ a la lista de primos.

\subsection{Criba Segmentada}
En \cite{Sorenson90} Sorenson dice que quizás la más práctica mejora de la criba de Eratóstenes sea la idea de segmentación. \\
En vez de cribar todo el intérvalo [2,$N$], se criban intérvalos de tamaño $\Delta$ uno a la vez. Se termina de hallar todos los números primos luego de cribar $\lceil(N-\sqrt{N})/\Delta\rceil$ intérvalos, siendo el último de longitud menor o igual a $\Delta$. La idea consta de dos pasos: \\
Primero, hallar todos los primos hasta $\sqrt{N}$.\\

Segundo, con $L = \sqrt{N}$, cribar el siguiente segmento [$L+1,L+\Delta$]; de los subsiguientes segmentos, su límite izquierdo se obtiene recurrentemente en base a $L = L+\Delta$.\\ 

Cuando $N$ es grande y no hay suficiente espacio para tener toda la lista $[2,N]$ en memoria, pero podemos manejar sublistas de tamaño $\Delta$, esta segmentación en una buena salida (obviamente sólo funciona si la lista hasta $\sqrt{N}$ es manejable, dados los límites en los recursos).

\begin{ejemplo} Con $N=27$. En vez de cribar de una vez todo el rango: \\
2, 3, \cancel{4}, 5, \cancel{6}, 7, \cancel{8}, \cancel{9}, \cancel{10}, 11, \cancel{12}, 13, \cancel{14}, \cancel{15}, \cancel{16}, 17, \cancel{18}, 19, \cancel{20}, \cancel{21}, \cancel{22}, 23, \cancel{24}, \cancel{25}, \cancel{26}, \cancel{27}\\
Elegimos un tamaño de segmento $\Delta = 9$,\\
$L = \sqrt{N} = \sqrt{27} = 5.19615$. \\
Primer paso: Obviamente trabajamos con la parte entera, es decir, hallamos los primos hasta $\lfloor\sqrt{27}\rfloor = 5$ de algún modo, una alternativa es la criba clásica de Eratóstenes.\\
En el ejemplo, luego de tachar 4, la lista inicial o base de primos es [2,3,5].\\
Segundo paso:\\
Segmento [$L+1,L+\Delta$] = [6,14].\\
Ya explicamos en el Enunciado~\ref{enun:elimPrevia} que el primer número compuesto útil del primo $p$ es $p^{2}$.
El primer múltiplo de 2 en este segmento no es $2^{2}$ sino 6, y se tachan 6, 8, 10, 12 y 14.\\
El primer múltiplo útil de 3 en este segmento es $3^{2}$, y se tachan 9 y 12.\\
Primos del primer segmento: 7, 11, 13.\\
Segmento [15,23].\\
El primer primo útil de 2 es 16, tachamos 16, 18, 20, 22.\\
El primer primo útil de 3 es 15, tachamos 15, 18, 21.\\
Primos del segundo segmento: 17, 19, 23.\\
Primos de los segmentos: 7, 11, 13, 17, 19, 23.\\
Segmento [24,27].\\
El primer primo útil de 2 es 24, tachamos 24, 26.\\
El primer primo útil de 3 es 24, tachamos 24, 27.\\
El primer primo útil de 5 es 25, tachamos 25.\\
No hay primos del tercer segmento.\\
Primos de los segmentos: 7, 11, 13, 17, 19, 23.\\
Primos (base y de los segmentos): [2, 3, 5, 7, 11, 13, 17, 19, 23].
\end{ejemplo}
El ejemplo tiene varias aspectos que hay que trabajar.\\
Ignoremos por un momento la optimización de que para cada primo base $p$ es mejor empezar a tachar los múltiplo desde $p^{2}$.\\
De modo que para cada primo base $p$ -en orden creciente- debemos calcular el primer múltiplo de dicho primo en cada nuevo segmento.

\begin{ejemplo}
Sea $N=100$ (en vez de 27). Sea $p=7$. Sea $\Delta$=10\\
Inicialmente $L = \sqrt{100} = 10$, pero con las actualizaciones llega a tomar el valor $L = 50$.\\
Segmento [51,60] ¿Cuál es el primer múltiplo de 7 en este segmento?\\
La respuesta es 56 = 50+7-(50 mod 7). A partir de ahí tachamos los siguientes múltiplos de 7 añadiendo 7 como es usual.
\end{ejemplo}

De modo que es crucial, dado el segmento [$L+1,L+\Delta$], hallar el primer múltiplo de $p$ en este segmento. Lo demás (tachar) es algo que ya conocemos. La expresión $L+p - (L \mod p)$ obtiene dicho primer múltiplo.

\begin{enunciado}[Primer múltiplo en un segmento]\label{enun:1ermultiplosegmento}\leavevmode\par
Sea $N \in \Z$, el número hasta el cual queremos hallar la lista de primos.\\
Sea $L \geq \lfloor \sqrt{N} \rfloor$.\\
Sea el número primo base $p \leq \lfloor \sqrt{N} \rfloor$.\\ 
Sea $L+1$ el límite inferior de un segmento de enteros de longitud $\Delta$.\\
Entonces $L+p - (L \; mod \; p)$ es el primer múltiplo de $p$ estrictamente mayor que $L$ y pertenece al segmento [$L+1,L+\Delta$] si y sólo si $p-(L \mod p) \leq \Delta$.
\end{enunciado}

\begin{proof} Por el algoritmo de la división, Enunciado~\ref{enun:algDivision}, $L=qp+r$, con $0 \leq r<p$.\\
a) Por definición de la operación módulo, $r=(L \mod p)$.\\
$L+p-(L \mod p)$ \\
$= (qp+r)+p-r$ \\
$=qp+p$ \\
$=(q+1)p$ \\
Como (q+1) es un entero, ello muestra que $L+p-(L \mod p)$ es un múltiplo de $p$.\\
b) El múltiplo de $p$ inmediatamente anterior a $(q+1)p$ es $(q+1)p-p =qp+p-p=qp$.\\
Sabemos que $0 \leq r$ y por tanto $qp \leq qp+r$, es decir,  $qp \leq L$.\\
Se concluye que el múltiplo de $p$ inmediatamente anterior a $L+p-(L \mod p)$ no está en el segmento de interés.\\
c) Como $r<p$, entonces $p-r>0$ y por tanto $L+p-r>L$. Es decir, $L+p-(L \; mod \; p)>L$. Es decir, $L+p-(L \mod p) \geq L+1$.\\
Los últimos dos incisos muestran que $L+p - (L \mod p)$ es el primer múltiplo de $p$ mayor o igual a $L+1$.\\
En estricto rigor, para que $L+p-(L \mod p)$ esté en el segmento [$L+1,L+\Delta$], es preciso que  $L+p-(L \mod p) \leq L+\Delta$, es decir, que $p-r \leq \Delta$.
\end{proof}

En la práctica, cuando los primos base $p$ son grandes es frecuente que $p-(L \mod p) > \Delta$. En estos casos, el primo $p$ no tiene representantes en el segmento actual. \\
Es un comportamiento esperado y correcto del algoritmo: el flujo de control debe ignorar automáticamente a $p$ para ese segmento, garantizando que solo se procesen los primos cuyos múltiplos pertenecen al segmento actual.\\
Eso es muy notorio en el último segmento cuando éste se trunca cerca de N.

\begin{ejemplo}
Sea $N=34$. Sea $p=5$. Sea $\Delta$=10.\\
$\lfloor \sqrt{N} \rfloor = \lfloor \sqrt{34} \rfloor =  5$.
Los segmentos son [11,20], [21,30], [31,34].\\
Con $L=30$, el primer múltiplo de $p=5$ en el segmento [31,34] es $L+p-(L \mod p)=30+5-(30 \mod 5) =  35$, fuera del segmento actual.
\end{ejemplo}

El Algoritmo~\ref{alg:AlgCribaSegment} implementa esta idea.\\
Note que para el segmento [$L+1,L+\Delta$] = [izq,der], se hace izq = $L+1$.\\

Lo usual es der = $L+\Delta$, pero como el límite a considerar es $N$, se efectúa der = $\min(L+\Delta,N)$, en particular para el último segmento.\\ 

Ya discutimos que el primer múltiplo útil del primo $p$ es $p^{2}$ para no tachar múltiplos ya tachados por otros primos menores.\\
Y que el primer múltiplo del primo $p$ en el segmento [$L+1,L+\Delta$] es $L+p - (L \mod p)$.\\
Se prefiere tomar como primer múltiplo útil:  $\max(p^{2}, L+p - (L \mod p))$; una u otra expresión evitará comparaciones duplicadas.\\
El costo de calcular el máximo es despreciable frente al beneficio de evitar todo el trabajo redundante.\\ 

Ya advertimos que puede haber casos en los que, para $p, \Delta y N$ dados, el primer múltiplo puede quedar fuera del segmento, como un comportamiento esperado. El algoritmo puede controlar ello de un modo semejante a este:\\
\texttt{for j desde primermultiploutil hasta limiteder step p}\\

Si el primer múltiplo sobrepasa la variable limiteder simplemente pasamos al siguiente primo $p$.

\begin{lstlisting}[caption={Criba Segmentada}, label={alg:AlgCribaSegment}, mathescape=true]
# Entrada: N, Delta
# Salida: Lista de primos entre 0 y N

# Paso 1: Criba clasica hasta raiz(N) 
L = $\lfloor \sqrt N \rfloor$
primos_iniciales = [True,$\dots$,True]  # Lista de L+1 elementos True
primos_iniciales[0] = primos_iniciales[1] = False
p = 2
while $p^{2} \leq$ L:
    if primos_iniciales[p]:
        for i desde $p^{2}$ hasta L step p:
            primos_iniciales[i] = False
    p = p+1
primos_base = []
for i desde 2 hasta L:
    if primos_iniciales[i] = True:
        agregar i a primos_base

# Paso 2: Cribar segmentos
resultado = primos_base  # primos base: parte del resultado
while L < N:
    izq = L+1
    der = min(L+Delta, N)    # N trunca el ultimo segmento
    marcado = [True,$\dots$,True]  # (der-izq+1) elementos True
    for p en primos_base:
        # hallar 1er multiplo de p en [izq,der]. 
        # Dado que L=izq-1
        # entonces (izq-1)+p-((izq-1) mod p) = L+p-(L mod p)
        1erMultiploUtil = max($p^{2}$, (izq-1)+p-((izq-1) mod p))
        for j desde 1erMultiploUtil hasta der step p:
            # marcado[j] marca la posicion j en [izq,der]
            # ver siguiente comentario por el uso de [j-izq] 
            marcado[j-izq] = False  
    # marcado va de [izq,der] con i dentro, interesa marcado[i]
    # pero en los hechos (en el algoritmo) marcado empieza en 0 
    # entonces reindexar por desplazamiento
    for i desde izq hasta der:
        if marcado[i-izq]:     
            agregar i a resultado
    L = L+Delta
print(resultado)
\end{lstlisting}

En el algoritmo presentado se codifica:\\
\texttt{1erMultiploUtil = max($p^2$, (izq-1)+p-((izq-1) mod p))}\\
Dado que $izq=L+1$, es equivalente a codificar:\\
\texttt{1erMultiploUtil = max($p^2$, L+p-(L mod p))}\\
Sin embargo, es muy usual en vez de dicha instrucción tener esta:\\
\texttt{1erMultiploUtil = max($p^2$, ((L+p)//p)$\cdot$p)}\\
Dado que $L=izq-1$, se codifica como\\
\texttt{1erMultiploUtil = max($p^2$, ((izq-1+p)//p)$\cdot$p) = max($p^2$, $\lfloor \frac{izq-1+p}{p} \rfloor \cdot p$)}.\\
El motivo se desarrolla a continuación.\\

Por el algoritmo de la división, Enunciado~\ref{enun:algDivision}, dividir $L+p$ entre $p$ devuelve $L + p = qp + r$, con $0 \leq r < p$.\\
Despejando $qp$ se tiene $qp = L + p - r$.\\
Dado que $(p \mod p) = 0$, es obvio que $r = (L + p) \mod p = L \mod p$, por tanto:\\
$L + p - (L \mod p) = qp = \lfloor \frac{(L + p)}{p} \rfloor p$, pues el cociente es $q = \lfloor \frac{L+p}{p} \rfloor$.

Y como en algunos lenguajes de programación la parte entera de la división se obtiene con $"//"$, resulta que $L + p - (L \mod p) = ((L + p) // p) p$.
 
$((L + p) // p) p$ es preferible a $L+p-(L \mod p)$ pues la operación módulo es más costosa que la división entera, además la expresión elegida es más usual en programación competitiva, es decir, coincide con el patrón que casi todo el mundo usa para ese problema, siendo fácil de reconocer y minimiza errores.

Es muy usual simular la función techo (ceil) con la función parte entera (floor).\\
En programación competitiva, leer código rápido es casi tan importante como escribirlo rápido. Al ver $((n + k - 1) // k) k$, se entiende que: “se está redondeando $n$ hacia arriba al siguiente múltiplo de $k$" (obtiene “el primer múltiplo de $k$ mayor o igual a n"). \\
La expresión es clásica porque es la forma más eficiente y elegante de realizar un redondeo hacia arriba (ceiling division) al múltiplo más cercano de $k$.

Lógica:\\
- $n // k$ trunca hacia abajo.\\
- Al sumar $k - 1$, si hay resto, se fuerza a subir al siguiente entero.\\
- Si $n$ ya es múltiplo de $k$: Sumar $k-1$ no es suficiente para alcanzar el siguiente múltiplo, por lo que la división se mantiene igual.\\
- Si $n$ tiene un resto (aunque sea 1), sumar $k-1$ empuja el resultado justo por encima del límite del siguiente múltiplo, haciendo que la división entera aumente en 1.\\
- Multiplicar por $k$ da el múltiplo exacto superior.\\
En programación competitiva, donde la velocidad de ejecución y la brevedad del código importan, esta fórmula evita usar funciones de punto flotante como ceil() (por ejemplo de la librería math), las cuales pueden ser lentas o tener problemas de precisión con números muy grandes. Y es idiomática porque aparece constantemente en paginación, alineación de memoria o división en bloques.\\
 
En efecto, $\lfloor \frac{izq-1+p}{p} \rfloor = \lceil izq/p \rceil$.\\
En general, por el Enunciado~\ref{enun:ceilequivalentefloor},  $\lfloor \frac{a-1+b}{b} \rfloor = \lceil a/b \rceil$.\\
Como $izq=L+1$:\\
$\lfloor \frac{izq-1+p}{p} \rfloor = \lceil \frac{L+1}{p} \rceil$.\\
Ello nos sirve para probar lo siguiente.

\begin{enunciado}[Primer múltiplo como Techo (ceil)]\label{enun:1ermultiploTecho}\leavevmode\par
$\lceil \frac{L+1}{p} \rceil p = L+p - (L \mod p)$
\end{enunciado}

\begin{proof}
Por el algoritmo de la división: $L = qp + r$ con $0 \leq r < p$ y por tanto con $1 \leq r+1 \leq p$.\\
Caso $r=0$: \\
$\lceil \frac{L+1}{p} \rceil = \lceil \frac{qp+1}{p} \rceil$ \; // pues $L = qp+r$ y $r=0$\\
$= \lceil q + \frac{1}{p} \rceil$ \; \; \; \; \; // obvio\\
$= q+1$\\
Por tanto:\\
$(q+1)p = qp+p$ \; \; \; \; \; \; \; // distribuyendo\\
$= L+p$ \; \; \; \; \; \; \; \; \; \; \; \;  \; \; \; // pues $L = qp+r$ y $r=0$\\
$= L+p - 0$  \; \; \; \; \; \; \; \; \;  \; \; \; // obvio\\
$= L+p - r$ \; \;  \; \; \; \; \; \; \; \; \; \; \;// pues $r=0$\\
$= L+p - (L \mod p)$\; \; \; \; // pues $r = (L \mod p)$\\
Caso $r>0$ y $r+1 < p$: \\
$\lceil \frac{L+1}{p} \rceil = \lceil \frac{qp+r+1}{p} \rceil$ \; \; \; \; \; \; \; \;// pues $L = qp+r$ \\
$= \lceil q + \frac{r+1}{p} \rceil$  \; \; \; \;  \; \; \; \;  \; \; \; \; // obvio  \\
$= q+1$  \; \; \; \;  \; \; \; \;  \; \; \; \;  \; \; \; \;// pues $r+1 < p$  \\
Por tanto: \\
$(q+1)p = qp+p$  \; \; \; \;  \; \; \; // distribuyendo \\
$= qp+r+p-r$  \; \; \; \;  \; \; \; \;  \;  // obvio  \\

$= L+p-r$  \; \; \; \;  \; \; \; \;  \; \; \; \;  \; \; \; \; // pues $L = qp+r$  \\
$= L+p - (L \mod p)$  \; \; \; \;  \; \; \; \; // pues $r = (L \mod p)$  \\
Caso $r>0$ y $r+1=p$: \\
Como $r+1=p$, entonces $r=p-1$ y también $1=p-r$.\\
$\lceil \frac{L+1}{p} \rceil = \lceil \frac{qp+r+1}{p} \rceil$  \; \; \; \;  \; \; \; \;  \; \; \; \; // pues $L = qp+r$ \\
$= \lceil q + \frac{r+1}{p} \rceil$  \; \; \; \;  \; \; \; \;  \; \; \; \;  \; \; \; \; \; // obvio \\
$= \lceil q + \frac{p}{p} \rceil$  \; \; \; \;  \; \; \; \;  \; \; \; \;  \; \; \; \;  \; \;  // pues $r+1=p$  \\
$= q+1$  \; \; \; \;  \; \; \; \;  \; \; \; \;  \; \; \; \;  \; \; \; \; // obvio \\
Por tanto:\\
$(q+1)p = qp+p$  \; \; \; \;  \; \; \; \;  \; \; \; \; // distribuyendo  \\
$= qp+r+p-r$  \; \; \; \;  \; \; \; \;  \; \; \; \;  \; // obvio  \\
$= L+p-r$  \; \; \; \;  \; \; \; \;  \; \; \; \;  \; \; \; \; \; // pues $L = qp+r$  \\
$= L+1$ \; \; \; \;  \; \; \; \;  \; \; \; \;  \; \; \; \;  \; \; \; \;  // pues $1=p-r$  \\
$= L+ (p-(p-1))$  \; \; \; \;  \; \; \; \;  \; \; \; // obvio  \\
$= L+ (p-r)$  \; \; \; \;  \; \; \; \;  \; \; \; \;  \; \; \; \; // pues $r=p-1$  \\
$= L+p - (L \mod p)$  \; \; \; \;  \; \; \; \;  \; // pues $r = (L \mod p)$  \\
En todos los casos tenemos que $\lceil \frac{L+1}{p} \rceil p = L+p - (L \mod p)$.
\end{proof}

\subsection{Criba de Sundaram}
Se adjudica la autoría de esta criba al estudiante de la India S. P. Sundaram.\\
La idea central y su justificación son simples y al mismo tiempo novedosas.
Es semejante a la criba de Eratóstenes, pero trabaja sólo con números impares y detecta sólo los primos impares.
Todos los números compuestos impares, pueden pensarse como el producto de dos número impares.

\begin{enunciado}[Múltiplos impares como producto de dos impares]\label{enun:multiploimpar}\leavevmode\par
Si $n$ es compuesto e impar entonces $n= a \cdot b$, con $a,b > 1$ e impares.
\end{enunciado}

\begin{proof}
Sea $n$ impar y compuesto.
Por ser compuesto $n$ tiene al menos un divisor propio $a>1$. Luego $n = a \cdot b$, con $b > 1$. \\
Si $a$ es par $a \cdot b$ es par.
Si $b$ es par $a \cdot b$ es par.
Si $a,b$ son pares $a \cdot b$ es par.\\
Como $n$ es impar se deduce que $a$ y $b$ son impares.
\end{proof}

Todo número impar mayor a 1 puede escribirse como $2k+1$, con $k>0$. Luego: \\
$n = a \cdot b = (2i+1)\cdot(2j+1) = 2(i+j+2ij)+1$.\\
Sea $m = i+j+2ij$. Entonces todo número impar compuesto $n$ puede escribirse en la forma $n = 2m+1$, con $m = i+j+2ij$. \\
La criba de Sundaram elimina estos valores de $m$.\\
Los $m$ que no son eliminados no corresponden a ningún producto de dos números impares mayores que 1. \\
Con los valores $m$ que quedan, los números $2m+1$ son primos impares.\\
Dado un límite $N$, se generan los números primos impares hasta $2N+1$.\\
El detalle está en generar los valores $m = i+j+2ij$. \\
En el ejemplo lo haremos de un modo más manual.\\
La idea puede implentarse con conjuntos, listas o un bitarray.

\begin{ejemplo}
Con $N=15$. \\  
Generar todos los pares $(i,j)$ con $(i+j+2ij) \leq 15$.\\
$i=1: \\
. j=1 \Rightarrow 1+1+2=4   \\
. j=2 \Rightarrow 1+2+4=7   \\ 
. j=3 \Rightarrow 1+3+6=10   \\
. j=4 \Rightarrow 1+4+8=13   \\ 
. j=5 \Rightarrow 1+5+10=16 > 15$  (nos detenemosm)  \\
$i=2:  
. j=2 \Rightarrow 2+2+8=12   \\
. j=3 \Rightarrow 2+3+12=17 > 15 \\ 
i=3:  \\
. j=3 \Rightarrow 3+3+18=24 > 15$ \\
No tiene sentido probar con valores mayores. \\
Por tanto, los números eliminados son: \{4,7,10,12,13\}. \\
Lo números que quedan hasta $N$ son: \{1,2,3,5,6,8,9,11,14,15\}.\\
Obtenemos $2m+1$ con cada uno: \{3,5,7,11,13,17,19,23,29,31\}, son los primos impares hasta $2N+1 = 31$.\\
Obviamente hay que añadir el primo par $2$.
\end{ejemplo}

La idea se implementa en el Algoritmo~\ref{alg:AlgCribaSundaram}.

\begin{lstlisting}[caption={Criba de Sundaram}, label={alg:AlgCribaSundaram}, mathescape=true]
# Entrada: N
# Salida: Lista de primos impares entre 1 y 2N+1

eliminados = [ ]
# generar numeros de la forma i + j + 2ij
i = 1
while $i \leq N$:
    j = i
    while  $i + j + 2ij \leq N$:
         agregar  $i + j + 2ij$ a eliminados
         j = j+1
    i = i+1
# construir lista de los que quedan
restantes = [ ]
for k=1 hasta N:
    if k not in eliminados:
        agregar k a restantes
# obtener $2m+1$
primos = [ ]
agregar 2 a primos  
for r in restantes:
    agregar (2r + 1) a primos
print(primos)
\end{lstlisting}

No compite en eficiencia pero es muy elegante.

\subsection{Criba de Atkin}
Como se vio, hay variantes y otras ideas sobre las cribas.\\
Aunque algunas de ellas tienen complejidades teóricas mejores, en la práctica suelen ser más lentas y difíciles de implementar, por las estructuras que requieren. \\ \\
Las cribas sublineales, como la de Atkin, requieren
operaciones internas que ocultan constantes en su complejidad temporal y que, en software real, terminan reduciendo su rendimiento. De este modo, pierden ventaja frente a, por ejemplo,  optimizaciones de la criba de Eratóstenes. \\
Sin embargo, muestran otros enfoques así como elegantes demostraciones teóricas que trascienden la literatura académica. Además, podrían beneficiarse de futuros avances tecnológicos que las adapten mejor, otorgándoles eventualmente un lugar privilegiado en la práctica. \\

La criba de Atkin \cite{AtkinBernstein2004} es un algoritmo moderno para generar todos los números primos hasta un límite dado $N$. Fue propuesta por Atkin y Bernstein y se basa en propiedades de las formas cuadráticas y en el uso del módulo 60 (wheel 60), que elimina de antemano los múltiplos de los primos pequeños 2, 3 y 5. \\

En el Anexo~\ref{anexo:iso} se presentan algunas pautas de la demostración de Atkin que fundamenta su criba. En las siguientes líneas presentaremos algunos conceptos e ideas que describen de manera simplificada parte de esas pautas. \\
Primero se introduce lo que se denomina forma cuadrática y la idea de que un entero puede representarse mediante formas cuadráticas específicas, como solución de ellas. \\

$Q(x,y) = ax^2 + bxy + cy^2$, con $a,b,c$ coeficientes
enteros, es una forma cuadrática genérica en dos variables. \\
Una más específica para $a=1, b=0, c=1$ es: $x^2+y^2$. \\

Cualquier entero $n$ puede representarse como la solución a la
ecuación con forma cuadrática $ax^{2} + bxy + cy^{2} = n$, con
$a = n,b = 0,c = 0$; $(x,y) = (1,0)$ es solución. \\
Pero no interesa este artificio sino formas cuadráticas específicas como $x^{2} + y^{2}$ o ${4x}^{2} + y^{2}$, para ver qué enteros admite tal o cual forma cuadrática. \\

\begin{ejemplo}
El entero 7 no es una suma de dos cuadrados, no puede representarse como
solución de la forma cuadrática $x^{2} + y^{2}$. \\
El 7 sí puede representarse como $x^{2} + 3y^2$ con
(x,y) = (2,1). \\
$4x^2 + y^{2}$ admite el entero 13 pues $4x^2 + y^2 = 13$
con $(x,y) = (1,3)$. \\
El 13 puede representarse como solución de la forma cuadrática
$4x^2 + y^2$. Esta solución con $x,y$ enteros positivos es
única (puede verificarse manualmente o verificarse con un programa, o verificarse en https://www.wolframalpha.com con la consulta "$4x^2 + y^2 = 13$, $x > 0, y > 0$ integer solutions"). \\
$4x^2 + y^2$ admite el entero 65 pues $4x^2 + y^2 = 65$ con $(x,y) = (4,1)$. \\
También $4x^2 + y^2 = 65$ con $(x,y) = (2,7)$. Estas dos
soluciones con $x,y$ enteros positivos son únicas. Luego, el número de soluciones con enteros positivos es par en este caso. \\
$4x^2 + y^2$ admite el entero 73 pues $4x^2 + y^2 = 73$ con $(x,y) = (4,3)$. Esta solución es única. Luego, el número de
soluciones con enteros positivos es impar en este caso. \\
$4x^2 + y^2$ admite el entero 325 con $(x,y) = (3,17),
(x,y) = (5,15)$ y $(x,y) = (9,1)$.\\
Estas soluciones son únicas. Luego, el número de soluciones con enteros positivos es impar en este caso.
\end{ejemplo}

Sea $n$ el entero del cual queremos decidir si la forma cuadrática $4x^2 + y^2 = n$ lo admite y cuántas soluciones postivas hay en ese caso. \\
Como $y^2 \geq 0$, se deduce que $4x^2 \le n$, luego: $x < \sqrt{\frac{n}{4}}$. \\ 

Cuando se dice que la verificación del número de soluciones y su paridad puede hacerse por programa, hablamos del siguiente algoritmo.

\begin{lstlisting}[caption={Verificación del número de soluciones enteras positivas}, label={alg:AlgVerifica}, mathescape=true]
num_sol = 0
for x=1 hasta $\lfloor \sqrt{\frac{n}{4}} \rfloor$:
    $y_2$ = $n - 4x^2$ 
    if $y_2$ > 0:
        y = $\lfloor \sqrt{y_2} \rfloor$
        if $y^2$ = $y_2$:        # testea si y es solucion exacta
            num_ sol = num_ sol + 1
\end{lstlisting}

Dado el número de soluciones num\_sol, se puede determinar si es par o no. \\

Atkin y Bernstein, vieron que la paridad del número de representaciones (o soluciones enteras) positivas de una forma cuadrática específica para un número $n$ dado contiene información sobre su posible primalidad. \\

Otra herramienta que usan es el módulo $M = 60 =
2^2 \cdot 3 \cdot 5$. \\
CRR(60) = \{1,7,11,13,17,19,23,29,31,37,41,43,47,49,53,59\}. \\
Recordando lo dicho en la criba por rueda, un número $n > 5$
coprimo con $M$ es candidato a primo. En cambio si $n > 5$ es
tal que $mcd(n,M) > 1$, es compuesto, en particular múltiplo de 2,
3 o 5. \\

También hacen uso del concepto de número libre de cuadrados. \\
Los números libres de cuadrados son números enteros que no pueden
dividirse por cuadrados perfectos distintos de 1, no son divisibles por $p^2$ para ningún primo $p$. \\
Es decir, su factorización principal tiene exactamente un factor para cada número primo que aparece en ella.

\begin{ejemplo}
Consideremos los siguientes casos: \\
$10 = 2 \cdot 5$ es un número libre de cuadrados. \\
$73$ es primo y un número libre de cuadrados. \\
$18 = 2 \cdot 3^2$ no es libre de cuadrados, $3^2$ divide a 18. \\
$45 = 3^2 \cdot 5$ no es libre de cuadrados, $3^2$ divide a 45. \\
$49 = 7^2$ no es libre de cuadrados, $7^2$ divide a 49. \\
\end{ejemplo}

Sea $n$ un número libre de cuadrados y $n \equiv 1 \pmod 4$. \\
En la implementación práctica, Atkin y Bernstein añaden la restricción de que $n$ tenga resto módulo 60 en CRR(60). Y subdividen CRR(60) en tres subconjuntos para trabajar cada uno de ellos con una forma cuadrática diferente. Usaremos ello con libertad. \\
Para la forma cuadrática $4x^2+y^2 = n$ usan \{1, 13, 17, 29, 37, 41, 49, 53\}. Es fácil probar que si $n$ es congruente con cualquiera de esos residuos módulo 60, también es congruente con 1 módulo 4. Por ejemplo: \\
Si  $n \equiv 17 \pmod {60}$ entonces $n = 60k + 17$.\\
Se sabe que $60k = 4 \cdot 15k \equiv 0 \pmod 4$.
Se sabe que $17 \equiv 1 \pmod 4$. \\
Entonces: $n = 60k+17 \equiv 1 \pmod 4$. \\ 
Uno de los tres enunciados en el artículo de Atkin es este: $n$ es primo si y solo si el número de soluciones enteras positivas de la ecuación $4x^2+y^2 = n$ es impar. \\
Y un enunciado, que incluye las tres clases, comúnmente aceptado es este:
\begin{enumerate}[label=\roman*)]
\item Un número $n$ libre de cuadrados, con resto módulo 60 en \{1, 13, 17, 29, 37, 41, 49, 53\} es primo si y solo si el número de soluciones enteras positivas de la ecuación $4x^2+y^2=n$ es impar.
\item Un número $n$ libre de cuadrados, con resto módulo 60 en \{7, 19, 31, 43\} es primo si y solo si el número de soluciones enteras positivas de la ecuación $3x^2+y^2=n$ es impar.
\item Un número $n$ libre de cuadrados, con resto módulo 60 en \{11, 23, 47, 59\} es primo si y solo si el número de soluciones enteras positivas de la ecuación $3x^2 - y^2=n$, con la condición $x>y>0$, es impar.
\end{enumerate}

Nos centraremos en el primer caso (el resto de casos se deja para el Anexo; aunque allí no se los trabaja exhaustivamente, sólo se deja la sensación de por dónde van las cosas que sustentan esta criba).

\begin{ejemplo}
Consideremos un caso de primalidad y otro de no primalidad. \\
$n = 61$. \\
$n \equiv 1 \pmod {60}$. \\
Puede usar el algoritmo presentado para verificar que $4x^2+y^2 = 61$ tiene un única solución positiva $(x,y) = (3,5)$. \\
El número de soluciones es impar y, en efecto, 61 es primo. Y es libre de cuadrados. \\
$n = 221$. \\
$n \equiv 41 \pmod {60}$. \\
Puede usar el algoritmo presentado para verificar que $4x^2+y^2 = 221$ tiene dos soluciones positivas $(x,y) = (5,11)$ y $(x,y) = (7,5)$. \\
El número de soluciones es par y, en efecto, $221 = 13 \cdot 17$ es compuesto. Y es libre de cuadrados. \\
\end{ejemplo}

Para tener una idea aproximada de por qué el primer enunciado es cierto requerimos varios conceptos de teoría de números. \\

\textbf{Extensiones de los enteros} \\
Se refiere a construir conjuntos más grandes que $\Z$ añadiendo
nuevas ``unidades básicas'' para operar.

\begin{ejemplo} Ponemos algunos casos. \\
\textbf{Enteros cuadráticos $\Z[\sqrt{d}]$}\\
Conjunto de números de la forma $a+b\sqrt{d}$ con $a,b
  \in \Z$.\\
  d es un entero que no es un cuadrado perfecto.\\
  $\Z[\sqrt{2}]$ es un caso. \\
\textbf{Enteros de Eisenstein $\Z[\omega]$}.\\
  Conjunto de números de la forma $a+b\omega$ con $a,b \in \Z$,
  donde $\omega = e^{2\pi i/3}$ es una raíz cúbica de la unidad.\\
\textbf{Enteros de Gauss $Z[i]$}\\
  Conjunto de números de la forma $a+bi$ con $a,b \in \Z$.\\
  Se obtiene añadiendo la unidad imaginaria $i$, donde $i^2 = -1$.
\end{ejemplo}

\textbf{Enteros de Gauss} \\
Nos interesan estos enteros. \\
$Z[i] = \{a+bi \, / \, a,b \in Z\}$. \\
Todo entero de Gauss $z = a+bi$ tiene lo que se denomina su conjugado: $\overline{z} =a-bi$. El conjugado de $3+4i$ es $3-4i$. \\
Hay una función particular en $Z[i]$ llamada norma, se define así: \\
$N(a+bi) = z \cdot \overline{z} = (a+bi) \cdot (a-bi) = a^2+b^2$.\\
Tiene la propiedad de multiplicatividad: $N(z_1z_2) =N(z_1)N(z_2)$. \\
$Z[i]$ es un anillo con las operaciones de suma y multiplicación (que es distributiva respecto a la suma). \\
En el Anexo se muestra que $Z[i]$ es un dominio euclídeo, es decir, que tiene definida una división con resto y un algoritmo para dicha división. \\
$\frac{7+i}{1+3i} = 1-2i$ pues $(1-2i)(1+3i) = 1+3i-2i+6 = 7+i$ \\
$\frac{7+2i}{2+i} = (3-i) + i$, cociente $3-i$, residuo $i$, pues $(3-i)(2+i) + i = (6+3i-2i+1) + i = 7+2i$. \\
La identidad multiplicativa (o elemento neutro de la multiplicación) en $\Z$ es el número 1. Para cualquier entero $x \in Z$, $x \cdot 1 = x$. \\
Se llama unidad a todo elemento de un anillo que tiene inverso multiplicativo dentro del mismo anillo. Es decir, $x$ es unidad si existe $y$ tal que $x \cdot y = 1$. \\
En $\Z$ sólo 1 y -1 son unidades.
En cambio en $Z[i]: 1,-1,i,-i$ son unidades (en efecto, $-i \cdot i = -i^2 = -(-1) = 1$, etc.). \\
En  $Z[i]$, dos enteros gaussianos son asociados si uno es el producto del otro por una unidad. \\
Un ideal $I$ en álgebra es un subconjunto no vacío dentro de un anillo R (por ring -no por reales-) que cumple: \\
. Ser cerrado bajo la resta: si $x,y \in I, x-y \in I$. \\
. Absorción de la multiplicación: si $x \in I, r \in R, rx \in I$. \\
Un ideal es algo así como un contenedor cerrado de múltiplos. \\
Un ideal principal es el que está formado únicamente por todos los múltiplos de un solo elemento. \\
En $\Z$ todos los ideales son principales pues están generados por su elemento positivo más pequeño (excluimos el caso 0). El ideal $I$ que consiste en todos los pares, es generado por el 2, es usual denotarlo por $2\Z$. \\
En $\Z[i]$ todos los ideales son principales pues están generados por el elemento con la norma más pequeña (excluimos el caso 0). El ideal $I$ que consiste en todos los múltiplos del generador $(1+i)$ se construye multiplicando todos los elementos de Z[i] por $(1+i)$, es usual denotarlo por $(1+i)\Z[i] = \{0,\pm(1+i), \pm(1-i),\pm2,\pm2i,\pm(2+2i), \dots \} =$ \{enteros gaussianos con norma par\}. \\
Ser un dominio euclídeo hace posible que todos los ideales sean principales, es decir, $Z[i]$ es un dominio de ideales principales (PID, Principal Ideal Domain); y de ahí un dominio de factorización única (DFU, Unique Factorization Domain). Esto es, generalizando el Teorema Fundamental de la Aritmética, todo elemento no nulo y no unidad se puede escribir como producto de elementos irreducibles (primos de Gauss). Esa factorización es única salvo por el orden y la multiplicación por unidades.\\
$65 = (1 + 2i)(1 - 2i)(2 + 3i)(2 - 3i)$. \\
Un primo de Gauss, es un entero de Gauss que no puede descomponerse en el producto de dos enteros gaussianos sin que uno de ellos sea una unidad, por ello se dice que son elementos irreducibles en $Z[i]$. \\
$65$ no es primo de Gauss, $7$ y $(1+2i)$ lo son. \\
Una forma rápida para identificar un primo de Gauss z es: \\
. Cuando $a,b \neq 0$, $N(z)$ es un primo racional (como se denomina también a los primos en $\Z$). \\
. Cuando $b = 0$ y $z=a \equiv 3 \pmod 4$. \\
$N(1+2i)=1^2+2^2=5$ y 5 es un primo racional. \\
$7 = 7+0i$ con $b=0$ y $7 \equiv 3 \pmod 4$. \\
Es interesante que los primos racionales $p$ se clasifiquen respecto al residuo que dejan módulo 4. \\
Se excluye $p \equiv 0 \pmod 4$ pues en ese caso $p$ no es primo sino múltiplo de 4. \\
Se excluye $p \equiv 2 \pmod 4$ pues en ese caso $p = 4k+2 = 2(k+1)$, es decir, $p$ es múltiplo de 2.  Salvo el caso particular de $p=2$ que es primo racional pero no es primo de Gauss pues $2 = (1+i)(1-i)$.\\
Los primos racionales $p$ tales que $p \equiv 1 \pmod 4$ siempre se factorizan en Z[i], no son primos de Gauss. $5 = (1+2i)(1-2i)$. \\
Los primos racionales $p$ tales que $p \equiv 3 \pmod 4$ quedan inertes, son primos de Gauss, sin cambios. En $\Z \; p = 3$, en $\Z[i] \; p = 3$. \\

Ahora sí consideremos el argumento de Atkin. \\
Note que en la ecuación $n = 4x^2+y^2 = (2x)^2+y^2 = (y + 2xi)(y - 2xi)$, $n$ es una suma de cuadrados por un lado y es el producto de dos enteros de Gauss conjugados, por otro. \\
Se considera el caso de un entero positivo $n$ cuyo módulo 60 esté en el subconjunto \{1, 13, 17, 29, 37, 41, 49, 53\} de $CRR(60)$ y por tanto $n \equiv 1 \pmod 4$ y que $n$ sea libre de cuadrados. \\
Recordemos que esta criba cuenta el número de soluciones positivas de $n = 4x^2+y^2$, cuando dicho número es impar se asegura que $n$ es primo. \\
Otra forma de ver ello es que, como $n = (y + 2xi)(y - 2xi)$, cuando el número de soluciones es impar (1 para ser precisos), $n$ se factoriza en Z[i] como producto de dos factores primos de Gauss conjugados. Si hubiera más de una solución positiva, los factores primos de Gauss de $n$ serían más que solo dos. El remate está en que se puede probar que si $n$ se factoriza en $\Z[i]$ exactamente como dos primos de Gauss conjugados, $n$ es un primo racional.

\begin{ejemplo} 
Consideremos un caso de primalidad y otro de no primalidad. \\
$n = 61$. \\
$4x^2+y^2 = 61$ tiene un única solución positiva $(x,y) = (3,5)$. 
El número de soluciones es impar (1 con precisión).
Reemplazando dicha solución en $61 = (y + 2xi)(y - 2xi)$, resulta $(5 + 6i)(5 - 6i) = 25 - 30i + 30i + 36 = 61$.
Como N(5 + 6i) = N(5 - 6i) = 61 y 61 es un primo racional, ambos factores son primos de Gauss. El primo racional $n = 61$ tiene exactamente dos factores primos de Gauss conjugados en su factorización en $\Z[i]$. \\
También $61 = (6 + 5i)(6 - 5i)$, pero ello no cambia en absoluto que $61$ se factorice exactamente como producto de dos primos de Gauss conjugados, es como si escribiéramos $61 = 61 \cdot 1$, en efecto, $(6 + 5i) = (5 - 6i) \cdot i$ y $(6 - 5i) = (5 + 6i) \cdot (-i)$, es sólo una multiplicación por unidades. \\
$n = 221$. \\
$4x^2+y^2 = 221$ tiene dos soluciones positivas $(x,y) = (5,11)$ y $(x,y) = (7,5)$. 
El número de soluciones es par (2 con precisión).
Reemplazando dichas soluciones en $221 = (y + 2xi)(y - 2xi)$, resulta $(11 + 10i)(11 - 10i) = 121 - 110i + 110i + 100 = 221$ y $(5 + 14i)(5 - 14i) = 25 - 70i + 70i + 196 = 221$. \\
Note que $N(11+10i)=11^2+10^2=221=13 \cdot 17$ no es un primo racional, es compuesto, luego $11+10i$ y $11-10i$ no son primos de Gauss.
Tampoco $(5 + 14i)$ ni $(5 - 14i)$. \\
En rigor, $221 = (2+3i)(2-3i)(1+4i)(1-4i)$, (como no es una unidad) es compuesto en $Z[i]$, con cada uno de esos factores siendo un primo de Gauss.\\
Note que $11 + 10i = (2+3i)(1+4i)(-i)$ y $5 + 14i = (2+3i)(1-4i)(i)$, multiplicar por una unidad está permitido. \\
Note que $13 \equiv 1 \pmod 4$.
\end{ejemplo}

En el Anexo mostraremos, con ayuda del teorema de Navidad (que también se prueba ahí), que $p \equiv 1 \pmod 4$ es un primo racional si y sólo si se factoriza en $\Z[i]$ exactamente en dos primos de Gauss conjugados. \\
Una de las partes de ello requiere mostrar que $p$ no es primo en $\Z[i]$, lo que se hace con ayuda de lo que llamaremos teorema de la raíz imaginaria modular: $a^2 \equiv -1 \pmod p$ cuando $p \equiv 1 \pmod 4$; que a su vez se prueba con ayuda del teorema de Wilson: $(p-1)! \equiv -1 \pmod p$. \\

El teorema de Navidad dice que un primo $p$ impar es una suma de cuadrados si y sólo si $p \equiv 1 \pmod 4$. \\

La prueba para este primer caso sigue, más o menos, el siguiente hilo conductor \\
Si se halla que la ecuación $4x^2+ y^2 = n$ tiene una sola solución con $x>0, y>0$, entonces, en $n =  (y + 2xi)(y - 2xi)$ sustituimos los valores de x,y. \\
Esta factorización gaussiana de $n$ con sólo dos factores primos de Gauss (uno y su conjugado), prueba que $n$ es primo. \\
Pues, razonando por reducción al absurdo, si suponemos que $n$ es compuesto en $\Z$, digamos $n=a \cdot b$, lleva inexorablemente a contradicciones: \\
$a,b$ -como caso esencial- son primos diferentes de 2, 3 y 5 y de sus múltiplos, por el $CRR(60)$. \\
. Si ambos $a,b$ son $\equiv 1 \pmod 4$, cada uno se descompone en dos primos gaussianos. Resultado: al menos cuatro primos gaussianos en la factorización de $n$. No dos. En consecuencia, por lo menos dos soluciones positivas. Contradice que sólo hay una. \\
. Si uno de los factores $a,b$ es $\equiv 3 \pmod 4$, ese factor permanece primo en $Z[i]$ . El otro, si es $\equiv 1 \pmod 4$, se descompone en dos. Resultado: al menos tres primos gaussianos en la factorización de $n$. No dos. En consecuencia, por lo menos dos soluciones positivas. Contradice que sólo hay una.\\
. Si ambos $a,b$ son $\equiv 3 \pmod 4$, entonces por el teorema de los dos cuadrados de Euler-Fermat, $n$ debe tener factores primos con exponente par. Contradice que $n$ es libre de cuadrados. \\

El teorema de los dos cuadrados de Euler-Fermat (caso general) dice que: Un entero $n$ es una suma de dos cuadrados si y sólo si en su factorización prima, todo primo $p \equiv 3 \pmod 4$ aparece con exponente par. \\
De los teoremas mencionados, se dan pistas o una prueba completa en el Anexo. \\
Conviene advertir que es un demostración plausible pero larga, para aportar claridad. La de Atkin es más corta pues usa algunos otros conceptos de álgebra avanzada para acortar la prueba y simplificarla en longitud, añadiendo algo más de densidad. \\

Estas ideas resultan en el Algoritmo~\ref{alg:AlgCribaAtkin}. \\
El vínculo con la factorización en $\Z[i]$ y los otros aspectos es la justificación teórica, no el mecanismo algorítmico directo; por ejemplo, en vez de contar el número de soluciones y luego verificar paridad para un $n$ dado, se calcula cuál $n$ corresponde a una forma cuadrática dada reemplazando valores de $x,y$ que se iteran y cada vez que $n$ aparece en una forma cuadrática válida, se invierte su estado, lo que refleja la paridad de soluciones en las ecuaciones. \\

\begin{lstlisting}[caption={Criba de Atkin}, label={alg:AlgCribaAtkin}, mathescape=true]
# Entrada: N  (se supone N>5)
# Salida: Lista de primos entre 1 y N

criba = [False, $\cdots$, False]  # Lista de N elementos False
primos = [ ]
for x=1 hasta $\lfloor \sqrt{N} \rfloor$:
    for y=1 hasta $\lfloor \sqrt{N} \rfloor$:
        $n = 4x^2 + y^2$
        if $n \leq N$ AND $n \mod 60 \in \{1,13,17,29,37,41,49,53\}$:
            criba[$n$] = NOT criba[$n$]   # inversion
        $n = 3x^2 + y^2$
        if $n \leq N$ AND $n \mod 60 \in \{7,19,31,43\}$:
            criba[$n$] = NOT criba[$n$]
        $n = 3x^2 - y^2$
        if $x > y$ AND $n \leq N$ AND $n \mod 60 \in \{11,23,47,59\}$:
            criba[$n$] = NOT criba[$n$]
# eliminar multiplos de cuadrados
for $n$=7 hasta $\lfloor \sqrt{N} \rfloor$:
    if criba[$n$] = True:
        k = $n^2$
        for multiplo=k hasta $N$ step k:
            criba[multiplo] = False
for $n$=7 hasta $N$:
    if criba[$n$] = True:
        Agregar $n$ a primos
print(primos)
\end{lstlisting}

\subsection{Comparación de las cribas}
Es necesario comparar las cribas para decidir en la práctica cual conviene implementar.

\begin{table}[h!]
\centering
\caption{Comparación de Algoritmos de Criba (C. abrevia Complejidad)}
\label{tab:cribas}
\begin{tabular}{|l|c|c|c|}
\hline
\textbf{Algoritmo} & \textbf{C. Temporal} & \textbf{C. Espacial} & \textbf{Uso Típico} \\
\hline
Eratóstenes Clásica & \(O(n \log \log n)\) & \(O(n)\) & General, pequeño y mediano \(n\) \\
\hline
Eratóstenes con Bits + Impares & \(O(n \log \log n)\) & \(O(n/2)\) & Optimización de memoria \\
\hline
Criba Incremental Impar & \(O(n \log \log n)\) & \(O(\sqrt{n})\) & Primos bajo demanda \\
\hline
Criba por Rueda & \(O(n \log \log n)\) & \(O(n / P)\) & \(n\) muy grande, optimizada \\
\hline
Criba de Gries y Misra & \(O(n \log \log n)\) & \(O(n)\) & Paralelización \\
\hline
Criba Segmentada & \(O(n \log \log n)\) & \(O(\sqrt{n})\) & \(n\) muy grande, memoria limitada \\
\hline
Criba de Sundaram & \(O(n \log n)\) & \(O(n)\) & Teórico / educativo \\
\hline
Criba de Atkin & \(O(n / \log \log n)\) & \(O(n)\) & Alto rendimiento, \(n\) gigante \\
\hline
\end{tabular}
\end{table}

En la práctica la criba de Eratóstenes es más eficiente y requiere menos memoria.